%=============================================================================
%  Programme brief for the faculty mentor -- MVP ASK.
%
%  This is the MINIMUM VIABLE variant: the smallest build that flies the
%  mission end to end and finds a person. It is not a competition entry, and
%  it says so. The tables, schedule and approval letters are generated by
%  tools/proposal/build_mvp_brief.py from hardware/bom/sourced_bom.py, the
%  same source the single- and three-aircraft briefs read, so none of the
%  three can disagree with another about a price.
%
%  Build: pdflatex mentor-brief-mvp x3
%=============================================================================
\documentclass[11pt,a4paper]{article}
\usepackage[T1]{fontenc}
\usepackage{mathptmx}
\usepackage[margin=19mm,top=22mm,bottom=22mm]{geometry}
% array must precede longtable: >{...} specifiers in a longtable preamble
% are otherwise undefined.
\usepackage{array}
\usepackage{booktabs}
\usepackage{longtable}
\usepackage{needspace}
\usepackage{microtype}
\usepackage{titlesec}
\usepackage[hidelinks]{hyperref}
\usepackage{fancyhdr}
\usepackage{graphicx}
% brainwave.pdf is the society mark, converted from SVG by
% tools/proposal/svg_logo_to_pdf.py -- pdflatex cannot read SVG, and this
% machine has no inkscape or ImageMagick.
\graphicspath{{images/}}

\newcommand{\rs}[1]{\mbox{INR~#1}}
% Table furniture, defined once so every table matches.
\newcommand{\thead}[1]{\textbf{#1}}
\setlength{\tabcolsep}{5pt}
\setlength{\aboverulesep}{0pt}
\setlength{\belowrulesep}{0pt}
\setlength{\extrarowheight}{2pt}
\renewcommand{\arraystretch}{1.12}

\titleformat{\section}{\large\bfseries}{\thesection.}{0.6em}{}
\titlespacing*{\section}{0pt}{12pt}{4pt}
\titleformat{\subsection}{\bfseries}{}{0pt}{}
\titlespacing*{\subsection}{0pt}{10pt}{3pt}
\setlength{\parskip}{0.45em}
\setlength{\parindent}{0pt}
\pagestyle{fancy}\fancyhf{}
\fancyhead[L]{\small Autonomous search and delivery --- minimum viable build}
\fancyhead[R]{\small\thepage}
\renewcommand{\headrulewidth}{0.4pt}

\begin{document}

\thispagestyle{empty}

% ---- letterhead ------------------------------------------------------------
\noindent
\raisebox{-0.5\height}{\includegraphics[height=16mm]{brainwave.pdf}}%
\hfill
\raisebox{-0.5\height}{\includegraphics[height=15mm]{thapar_logo.png}}

\vspace{3mm}\hrule\vspace{6mm}

\begin{center}
{\LARGE\bfseries An Autonomous UAV for\\[1.5mm]
Post-Disaster Search, Localisation and Payload Delivery}\\[2.5mm]
{\large Minimum viable build}
\end{center}

\vspace{1mm}\hrule\vspace{4mm}

%-----------------------------------------------------------------------------
\section{Executive summary}

Floods displace millions of people in India every year. The hard part of a
rescue is not reaching survivors; it is finding them.

We want to build one small aircraft, about a metre across and about 7\,kg, that
searches flooded ground on its own. One person sends it out. It flies itself,
spots people from the air, works out where those people are, and drops a relief
kit to each. Nothing depends on a phone network or the internet.

\textbf{This is the smallest version of that aircraft which still proves the
idea.} It flies the whole mission --- plan, search, find, locate, descend,
drop, return --- and nothing else. We have deliberately left out everything
that makes the finished system a competition entry rather than a working one:
the live video downlink, the RTK base that makes positioning
centimetre-accurate, the aircraft-to-aircraft radio. Section~\ref{sec:resources} lists every omission
and what it costs.
\textbf{It is not a competition entry, and we would rather say so now than
discover it later.}

\textbf{If it works, we would build it up to the full aircraft, and then build
two more and fly all three as a coordinated group.} The design is already a
three-aircraft design and the coordination software is already written and
tested in simulation. What we are asking for here is the first evidence.

The design and the software are done, and we have flown the whole mission in
simulation. Nothing has flown for real yet, so every figure here is calculated
rather than measured.

We would like to begin with \rs{3{,}733}: enough to test whether one motor
lifts what the design needs. After that the money comes in 12 small steps, each
released only once the step before it has produced a result we can show you.
Nothing over \rs{26{,}000} is spent until that first thrust measurement exists.
If every step is approved the whole build comes to \rs{1{,}98{,}670}.

%-----------------------------------------------------------------------------
\section{The problem}

Boats are slow, dangerous for the crew, and cannot see onto rooftops.
Helicopters are few and are needed for evacuation. Ordinary drones need a pilot
each, rely on infrastructure the flood has washed away, and cannot help anyone
once they have found them.

Finding people turns out to be the harder half. Someone in floodwater shows
only their head and shoulders --- about 40\,cm of them is visible from above.
Detection software works on a small picture of fixed size, so the usual
approach shrinks the photograph, and the person shrinks with it until there is
too little left to recognise. Using a smaller, cheaper model makes that worse,
not better.

There is also no public set of labelled flood photographs to learn from. Thermal
cameras do not rescue us either: someone sitting in water ends up close to the
temperature of the water.

%-----------------------------------------------------------------------------
\section{Our solution}

The operator draws a search area of any shape and the software plans a route
over it. We have tried this on thirteen awkward shapes and no route ever
wandered outside the area. After that the aircraft flies itself, and our tests
confirm nothing is commanded from the ground once a mission has begun. It spots
people and works out where they are on its own, carries four relief kits,
descends over a survivor, holds still, and drops one. If the radio fails, the
battery runs low, or the operator hits abort, it ends the same way: come home
and land.

The trick with the camera is not to shrink the picture. We cut each frame into
tiles and search them at full size, and we fly lower. Between them those two
changes make a person in the water nine times bigger in the image, which is
what lifts them clear of the range where this kind of software stops working.
Neither change is enough on its own. It costs four times as much computing,
which the accelerator we have chosen can supply, and the mission takes no
longer, because flying lower saves more climbing than the extra passes cost.

There is one thing we have not solved. We plan to look at each patch of ground
twice a second, which the accelerator manages comfortably. To be really
confident about a sighting we would want three looks a second, and that is half
as much work again --- more than we can count on. There may be a way round it:
most of what the camera sees is plainly open water, so a quick, cheap check
could throw those patches away and leave the expensive search to the rest. The
danger is that if the quick check throws away a patch with a person in it,
nothing tells us. No error, no warning, just someone never found. We will
measure how often that happens against public data before we rely on it.

\subsection{What this build can and cannot do}

This is worth being exact about, because it is the whole argument for the two
steps that would follow it.

The rate at which ground is searched is a property of the camera and the flying
speed, not of how many aircraft there are: about 1.4 hectares a minute. One
aircraft therefore sweeps \textbf{three and a half hectares} in the same three
minutes that three aircraft take to sweep ten. Asked to cover the full ten
hectares alone it needs about eight minutes of sweeping rather than three, and
a complete mission then runs to roughly thirteen minutes against a hover
endurance of about fifteen --- which leaves none of the reserve our own sizing
rule requires.

So this aircraft demonstrates the whole mission, end to end, over the
three-and-a-half-hectare share it would fly as one of three. Every part of the
system is exercised: survey planning, autonomous flight, detection,
localisation, descent, release, and return.

\textbf{Four things it cannot do}, each because of something we chose not to
buy. It cannot cover ten hectares within the time limit and still hold a
reserve --- that needs the other two aircraft. It cannot deliver more than the
four kits it carries. It cannot show the operator live video, because there is
no downlink; the ground station sees the aircraft's position, state and
reported findings over the same radio the safety pilot uses, which is what the
mission actually turns on. And it cannot prove \emph{how accurately} it locates
people, because the RTK base station that would measure that against surveyed
ground is left out: it serves a publication objective rather than a flight
requirement.

The last one deserves emphasis, because it is the one real loss. The aircraft
still locates survivors --- without a base its receiver works to about two
metres against a five-metre delivery requirement, better than twice the margin it
needs. What we cannot yet
publish is a \emph{measured} figure for how well it does so.

\subsection{One change to the design: the battery is bought, not built}

The design specifies an eighteen-cell pack assembled in-house. Building it
needs a spot welder and a skill we do not have, and eighteen cells hand-welded
in a student lab is the most hazardous thing in this whole build.

We therefore buy a ready-made 6S 16000\,mAh pack instead. The sizing model
states the requirement without reference to chemistry --- \emph{minimum
nameplate, any chemistry: 244\,Wh, 6S, at least 115\,A continuous} --- and the
bought pack gives about 365\,Wh and 400\,A, against a custom pack that would
have given 292\,Wh and about 110\,A. Its reserve margin rises from $+$20\,\%
to $+$27--36\,\%, depending on what the pack turns out to weigh. A smaller 12000\,mAh pack was considered and rejected: once
its own weight is counted it falls short of the reserve policy.
\textbf{What it costs is weight}: roughly 0.4--0.8\,kg more, which takes
thrust-to-weight from 2.0 to about 1.8--1.9 on the same motors unless the
thrust stand shows more thrust than the design assumed. Cycle life is also
shorter for this chemistry, so a second pack should be expected over a long
test campaign.

%-----------------------------------------------------------------------------
\section{Deliverables}

\begin{enumerate}\itemsep0pt
\item One aircraft flying the mission end to end under one operator.
\item A ground station that shows position, state and findings, and by design
cannot fly the aircraft. It is built for three video feeds and this build
provides none, so it runs on telemetry alone.
\item The autonomy software, with an automatic test suite behind it. The
multi-aircraft coordination is part of it and is already tested in simulation.
\item Measured detection accuracy and a verdict on the two-stage filter.
\item A publication: tiling, aerial person detection and low-power onboard
inference have not been combined for flood conditions.
\item Source code and results, published in either outcome.
\item A measured answer to the question of whether to build the rest.
\end{enumerate}

%-----------------------------------------------------------------------------
\section{Benefits}

% Five parallel items, so they are set as a list -- the same way Deliverables
% above sets its parallel items. As running paragraphs the labels did not hang
% and the block read as one mass of text.
\begin{itemize}\itemsep2pt
\item Three and a half hectares searched in about three minutes, and what comes
back is a set of coordinates a rescue team can act on, not a video feed
somebody has to sit and watch.
\item The whole thing costs about an hour of helicopter charter at Indian
rates, travels by road, and is run by two people.
\item A thrust stand, a ground station, a tested codebase and students who have
built all of it --- and all of that stays in the department afterwards.
\item Every part but one comes from an Indian supplier, so nothing sits at
customs and anything broken in testing is replaced within days.
\item The department finds out what a flying aircraft actually costs and
actually does before being asked to fund three of them --- for under a third of
what three would cost.
\end{itemize}

%-----------------------------------------------------------------------------
\section{Resources needed}\label{sec:resources}

Funding is requested in 13 small steps rather than as one grant. Each step buys
one named set of parts and is released only once the previous step has produced
a stated result. If a step fails, the programme stops there and nothing beyond
it has been spent.

The order follows engineering dependency: one motor is measured before the
other three are bought, and the aircraft is assembled, weighed and flown before
the ground segment is completed around it.

% No \clearpage here: the longtable defines \endhead, so it repeats its header
% on every continuation page and can start mid-page. Forcing a fresh page left
% the lower half of the preceding page blank.
\subsection{Phase schedule}

The third column is the amount released. The fourth is what must already be
true before it is released.

%============================================================================
% GENERATED FILE -- do not edit.
% Produced by tools/proposal/build_mvp_brief.py from
% hardware/bom/sourced_bom.py. Edit the BOM, not this file.
%============================================================================
\begingroup
\footnotesize
\setlength{\extrarowheight}{1.5pt}
\begin{longtable}{@{}>{\bfseries}r>{\raggedright\arraybackslash}p{7.3cm}r>{\raggedright\arraybackslash}p{5.6cm}@{}}
\toprule
\thead{\#} & \thead{Objective} & \thead{INR} & \thead{Released only after} \\
\midrule
\endfirsthead
\toprule
\thead{\#} & \thead{Objective} & \thead{INR} & \thead{Released only after} \\
\midrule
\endhead
\midrule
\multicolumn{2}{@{}r}{\thead{Total}} & \thead{1{,}96{,}225} & \\
\bottomrule
\endlastfoot
1 & Establish thrust and mass measurement capability & 3{,}733 & Approval \\
2 & Measure one motor, propeller and speed controller on the stand & 14{,}771 & Instruments commissioned \\
3 & Aircraft --- autopilot and control link & 12{,}626 & \textbf{Thrust verified} \\
4 & Aircraft --- onboard computer and AI accelerator & 31{,}402 & Autopilot bench-tested \\
5 & Aircraft --- satellite positioning and compass & 17{,}701 & Compute stack operating \\
6 & Aircraft --- camera and command link & 9{,}474 & Position fix acquired \\
7 & Aircraft --- airframe fabrication & 18{,}886 & Avionics integrated \\
8 & Aircraft --- battery pack and power distribution & 23{,}515 & Airframe fabricated \\
9 & Aircraft --- release servos and wiring & 3{,}075 & Pack bench-discharged \\
10 & Aircraft --- propulsion completed & 25{,}661 & Drive train installed \\
11 & Charging and manual override & 27{,}448 & Aircraft assembled and weighed \\
12 & Field consumables and assembly materials & 7{,}934 & Charging verified \\
\end{longtable}
\endgroup


Phase 3 commits the first significant expenditure and is released only on a
measured thrust figure. Nothing in this schedule replicates anything: every
line buys a part that goes into the one aircraft, an instrument that measures
it, or the ground station it reports to.

\subsection{Component selection}

The schedule gives the amount
\emph{released} per phase. The tables below give the parts each phase
\emph{buys}. A released amount is the parts cost plus tax where tax is still
owed, plus 15\,\% contingency. Across the programme that is \rs{1{,}71{,}313}
of parts against \rs{1{,}98{,}670} released.

For scale: building this up to the complete single aircraft --- adding back the
video downlink, the RTK base, the rangefinder and the coordination
radio --- would come to \rs{2{,}45{,}993}, and the full three-aircraft
programme to \rs{5{,}78{,}179}. Each step is asked for only on the evidence
produced by the one before it.

\textbf{Every line below but one is a current listing from a named Indian
supplier, priced and linked;} the consumables are a round allowance. The
schedule, the tables and the approval letters are all generated from that one
list, so they cannot drift apart --- and they are generated from the same list
as the single-aircraft and three-aircraft versions of this brief, so none of the
three can disagree with another about a price. The battery's maker does not
publish its mass, so it is weighed on arrival.

GST is charged only on lines quoted before tax; Indian retail prices already
include it. The full parts list, with every listing linked and the checks each part must pass before it is bought or flown, is submitted alongside this brief as a separate Bill of Materials.

\needspace{7\baselineskip}
Terms used below. \emph{RTK}: satellite positioning corrected against
a second receiver on a surveyed point, giving centimetre rather than metre
accuracy. \emph{6S}: six cells in series, giving the pack its 22.2\,V.
\emph{C}: discharge rate as a multiple of capacity --- 25\,C on a 12\,Ah pack
is 300\,A. \emph{kgf}: thrust written as the weight it would lift.
\emph{TOPS}: trillion operations per second, the accelerator's throughput.
\emph{ESC}: the electronic speed controller between battery and motor.
\emph{BMS}: the board that balances and protects the cells --- there is none
in this build, and the note after the tables says why.

%============================================================================
% GENERATED FILE -- do not edit.
% Produced by tools/proposal/build_mvp_brief.py from
% hardware/bom/sourced_bom.py. Edit the BOM, not this file.
%============================================================================
\subsubsection*{Measurement instruments --- phases 1--2 \quad \rs{3{,}246} of parts}

\begingroup\small
\begin{longtable}{@{}>{\raggedright\arraybackslash}p{2.6cm}>{\raggedright\arraybackslash}p{4.4cm}r>{\raggedright\arraybackslash}p{7.4cm}@{}}
\toprule
\thead{Item} & \thead{Model} & \thead{INR} & \thead{Why this part, and from whom} \\
\midrule
\endfirsthead
\toprule
\thead{Item} & \thead{Model} & \thead{INR} & \thead{Why this part, and from whom} \\
\midrule
\endhead
Load-cell amplifier & 7Semi HX711 24-bit ADC & 104 & Turns the load cell's signal into a reading we can log. \textit{Robu} \\
Load cell & REES52 20 kg load cell + HX711 breakout & 854 & The motors ship without a thrust curve, so we measure it ourselves. Ready-made stands cost around ten times this. \textit{Amazon India} \\
Thrust-stand mast & SURYAKUSH metal 2.1 m tripod & 373 & Holds the motor clear of the bench, so its own draught does not flatter the reading. \textit{Flipkart} \\
Fasteners & Hex socket-head assortment M3/M4/M5, SS304 & 1{,}140 & A carbon airframe is bolted, not glued. \textit{OnlyScrews} \\
Threadlocker & Loctite 243, 50 ml & 664 & Stops bolts working loose under motor vibration. Medium strength, so a joint can still be undone. \textit{Amazon India} \\
Throttle source & Servo tester / ESC consistency tester, PPM & 111 & A motor on a stand needs something to command it. The safety-pilot transmitter is not bought until phase 11, long after this test. \textit{Robokits} \\
\bottomrule
\end{longtable}
\endgroup

\subsubsection*{Per aircraft, avionics and sensing --- phases 3--6 \quad \rs{84{,}659} of parts}

\begingroup\small
\begin{longtable}{@{}>{\raggedright\arraybackslash}p{2.6cm}>{\raggedright\arraybackslash}p{4.4cm}r>{\raggedright\arraybackslash}p{7.4cm}@{}}
\toprule
\thead{Item} & \thead{Model} & \thead{INR} & \thead{Why this part, and from whom} \\
\midrule
\endfirsthead
\toprule
\thead{Item} & \thead{Model} & \thead{INR} & \thead{Why this part, and from whom} \\
\midrule
\endhead
Flight controller & Holybro Pixhawk 6C Mini & 22{,}049 & The autopilot: flies the aircraft and runs its safety behaviours. Previously a quotation, now ordinary retail. \textit{Robu} \\
FC vibration mount & Glass-fibre anti-vibration shock absorber & 161 & Isolates the autopilot from propeller vibration, which otherwise confuses its sensors. \textit{Robu} \\
Autopilot log card & SanDisk 32 GB microSDHC, class 10 & 1{,}549 & The autopilot ships without one and records nothing without it. This is a different slot on a different board from the computer's card above. \textit{Robocraze} \\
Companion computer & Raspberry Pi 5, 8 GB & 19{,}999 & The computer the autonomy runs on. The accelerator and the camera both plug into it, so neither works without it. \textit{Robu} \\
AI accelerator & Raspberry Pi AI HAT+, 26 TOPS & 11{,}749 & Runs the person-detection model fast enough to search at flying speed. \textit{Robu} \\
Compute cooling & Official Raspberry Pi 5 active cooler & 519 & The hottest moment is the setup on the ground, before the propellers are turning and cooling it. \textit{Robu} \\
Storage & SanDisk High Endurance 128 GB microSD & 3{,}989 & Records every frame and detection for a whole mission. Ordinary cards wear out at that rate. \textit{Flipkart} \\
Camera + lens & Arducam 12.3 MP 1/2.3 in HQ module, 6 mm CS lens & 6{,}799 & The search sensor, chosen so a person in water is large enough in the picture to be found. \textit{Robu} \\
Command receiver & ExpressLRS RX24T, 2.4 GHz & 1{,}439 & Carries the safety pilot's control and the aircraft's reporting on a single link, which saves a second radio. \textit{Robu} \\
Power module & Holybro PM07 & 5{,}249 & Reports battery voltage and current to the autopilot, which is what triggers a low-battery return. \textit{Robu} \\
BEC, primary & 8 A UBEC, 5/6 V out, 2-6S in & 3{,}939 & Powers the computer. Deliberately oversized: losing it in flight is the failure this line exists to prevent. \textit{Amazon India} \\
BEC, secondary & ReadytoSky 2-6S 5V/3A + 12V/3A switchable UBEC & 279 & A separate supply for the servos, so a jammed servo cannot disturb the autopilot. \textit{Robu} \\
GNSS receiver & Holybro Micro M9N GPS with case, NEO-M9N and compass & 6{,}939 & Positioning for aircraft 2 and 3, accurate to about two metres against a five metre requirement. Carries the compass the flight controller needs, and plugs straight into it. \textit{in-house} \\
\bottomrule
\end{longtable}
\endgroup

\subsubsection*{Per aircraft, airframe and drive --- phases 7--10 \quad \rs{65{,}152} of parts}

\begingroup\small
\begin{longtable}{@{}>{\raggedright\arraybackslash}p{2.6cm}>{\raggedright\arraybackslash}p{4.4cm}r>{\raggedright\arraybackslash}p{7.4cm}@{}}
\toprule
\thead{Item} & \thead{Model} & \thead{INR} & \thead{Why this part, and from whom} \\
\midrule
\endfirsthead
\toprule
\thead{Item} & \thead{Model} & \thead{INR} & \thead{Why this part, and from whom} \\
\midrule
\endhead
Motors & Tarot TL96020, 5008, 340 KV, 4 off & 18{,}812 & Twelve. The first is measured on the stand before the other eleven are bought. \textit{Robokits} \\
Propellers & 1855 counter-rotating carbon fibre, CW+CCW pair, 6 off & 9{,}846 & Sixteen for twelve positions. The item most often broken in flight testing, and the whole of our spares policy. \textit{Robokits} \\
Speed controllers & Darkmatter VISHNU 50 MK-III, 4-in-1, 50 A & 5{,}189 & One board drives all four motors. Made in India, and previously a quotation. \textit{Robu} \\
Arm tube & 3K carbon fibre, OD25 x ID23 x 1000 mm, 2 off & 4{,}518 & The four arms. Carbon, and comfortably strong --- the joints, not the tubes, are what limit the airframe. \textit{Robu} \\
Frame plate stock & 3K carbon fibre sheet, 300 x 200 x 2 mm & 2{,}895 & The arms are tubes; this is the plate they bolt to. One sheet yields the top and bottom of a centre section at this size. \textit{in-house} \\
Motor mounts & Tarot 25 mm motor mount, TL9602, 4 off & 4{,}476 & Joins motor to arm without machining the tube. This joint is the part that has to be right. \textit{Robu} \\
Arm clamps & T-bolt hose clamps 23-25 mm, SS304, 4 pcs & 3{,}012 & Bolted rather than bonded, so an arm damaged in a heavy landing is swapped in the field instead of scrapping the frame. \textit{Amazon India} \\
Landing gear & F450/F550 landing skid & 549 & Gives clearance underneath for the kit magazine, and a stance wide enough for an untidy landing. \textit{Robu} \\
Suspension springs & 5 mm stainless double torsion spring, 10 off & 20 & Let the legs absorb a hard landing rather than pass it into the arms and the battery. \textit{IndiaMART} \\
Printed parts & Pro-Range PETG-CF filament, 1.75 mm, 1 kg & 949 & Filament for the magazine, mounts and covers. Holds its shape in the sun, which cheaper filament does not. \textit{Robu} \\
Release servos & MG90S mini servo, 180 deg, 4 off & 596 & Open the four kit stations. Metal gears, because a plastic one can drop a kit without anyone knowing. \textit{Robocraze} \\
Pack fusing & ANL fuse holder + 150 A ANL fuses & 3{,}000 & Protects against a short circuit without cutting power at full throttle. \textit{Njour} \\
Pack retention & 300 mm LiPo strap, reusable, 2 off & 124 & Holds the battery down. It is a fifth of the aircraft's weight, so a strap that lets go is not a small problem. \textit{Robu} \\
Main leads & 10 AWG ultra-flexible silicone wire & 149 & The main power cable, sized for the highest current the motors ever draw rather than the average. \textit{Robu} \\
Power connectors & Amass XT90S anti-spark, M/F pair, 4 off & 720 & Anti-spark, so connecting a battery this size does not burn the contacts. \textit{Robokits} \\
Signal connectors & JST ZH 6-pin, 200 mm leads & 23 & Keyed, so the wiring can only go back together one way after a repair. \textit{Robu} \\
Antenna feeders & SMA M-F bulkhead RG316, 150 mm + adapters & 1{,}139 & Let the antennas mount on the airframe instead of hanging off the radios. \textit{Desertcart} \\
Insulation & PVC heat-shrink sleeve, 150 mm & 47 & Sleeving over the battery joints --- the one place on the aircraft where a rubbed wire is an immediate fire. \textit{Robu} \\
Battery pack & GenX 22.2V 6S 12000mAh 25C/50C LiPo, XT90 & 9{,}088 & 266 Wh against the 244 Wh the sizing model requires of any chemistry, and 25\,C (300\,A) against a 115\,A peak. Bought rather than built: the custom 18-cell pack it replaces needs a spot welder, and this is lighter, cheaper and has more current margin. \textit{Robokits} \\
\bottomrule
\end{longtable}
\endgroup

\subsubsection*{Flight-test support --- phases 11--12 \quad \rs{29{,}867} of parts}

\begingroup\small
\begin{longtable}{@{}>{\raggedright\arraybackslash}p{2.6cm}>{\raggedright\arraybackslash}p{4.4cm}r>{\raggedright\arraybackslash}p{7.4cm}@{}}
\toprule
\thead{Item} & \thead{Model} & \thead{INR} & \thead{Why this part, and from whom} \\
\midrule
\endfirsthead
\toprule
\thead{Item} & \thead{Model} & \thead{INR} & \thead{Why this part, and from whom} \\
\midrule
\endhead
Safety-pilot transmitter & RadioMaster TX12 MKII ELRS + RP1 receiver & 16{,}519 & Manual override for the safety pilot, independent of the autonomy. A rule requirement and a safety one. \textit{zBotic} \\
Battery charger & ToolKitRC M6D, 500 W, 25 A, dual & 7{,}349 & Two batteries recharged between flights is what sets how many tests fit into a day. \textit{Robu} \\
Cable management & Nylon zip ties, 350 mm, 100 pcs & 135 & Ties long enough to dress wiring around a one-metre airframe. \textit{Robu} \\
Heat-shrink kit & 328 pcs assorted 2:1 heat-shrink & 159 & Insulates the signal wiring during assembly. \textit{Robu} \\
Mounting tape & 3M VHB 5952, 6 mm x 8.2 m & 560 & Holds radios and receivers where a screw would need a fixing the printed part cannot carry. \textit{Amazon India} \\
Hook and loop & Self-adhesive hook and loop, 25 mm x 2 m & 145 & For the battery and anything else that must come out between flights without tools. \textit{Amazon India} \\
Consumables & Fasteners, tape, connectors, filament & 5{,}000 & Solder, abrasives, spare hardware, filament. A round figure, because itemising it would be false precision. \textit{in-house} \\
\bottomrule
\end{longtable}
\endgroup

\needspace{6\baselineskip}
\subsubsection*{Deliberately not bought}

\begingroup\small
\begin{longtable}{@{}>{\raggedright\arraybackslash}p{3.2cm}r>{\raggedright\arraybackslash}p{11.4cm}@{}}
\toprule
\thead{Item} & \thead{INR saved} & \thead{Why the MVP can fly without it} \\
\midrule
\endfirsthead
\toprule
\thead{Item} & \thead{INR saved} & \thead{Why the MVP can fly without it} \\
\midrule
\endhead
GNSS RTK receiver & 50{,}000 & research instrument. The delivery requirement is 5 m and the M9N gives about 2 m; RTK exists to MEASURE geolocation against surveyed ground, which is a publication objective, not a flight one \\
Pack health monitor & 9{,}058 & bench instrument for tracking pack ageing; does not fly \\
Video receivers & 4{,}399 & ground end of the same downlink \\
Video transmitter & 4{,}239 & FPV downlink. The mission is autonomous -- nothing has to watch it live \\
Height rangefinder & 2{,}827 & gates the release at 6 m over water of unknown level. The barometer suffices over dry ground, which is where the MVP flies \\
Coordination base & 1{,}646 & ground end of the same link. Telemetry reaches the laptop over ExpressLRS native MAVLink instead \\
Coordination radio & 989 & aircraft-to-aircraft link, and there is one aircraft \\
Base station mount & 879 & tripod for the RTK base that is no longer bought \\
Receive antennas & 839 & ground end of the same downlink \\
Video capture & 399 & ground end of the same downlink \\
\midrule
\multicolumn{1}{@{}r}{\thead{Parts not bought}} & \thead{75{,}275} & \\
\bottomrule
\end{longtable}
\endgroup

\needspace{6\baselineskip}
\textbf{Why there is no battery management board.} An earlier version funded
one. Most of what that line covered is already bought under other names, and
the part that remains --- a board that cuts power when a cell is unhappy --- is
the wrong device on an aircraft, because it can cut power at full throttle,
which is a crash rather than a protection. Instead the battery is fused against
short circuits, the autopilot brings the aircraft home on low voltage, and
every cell is checked on the charger before each flight. What we do not have is
per-cell protection in the air; that is a deliberate choice, not an oversight.

\subsection{What this arrangement costs us}

We should be honest that twelve approvals is a lot of paperwork for the
office, and that buying parts for one aircraft means single-unit prices with no
quantity break. If the paperwork becomes the bigger problem, the per-subsystem
steps can be merged into fewer orders.

We should also be honest that a cut is not always a saving. Everything in the
table above is a part the finished system needs; deferring it moves the cost
later and adds a second delivery, a second integration and a second round of
testing. What deferral buys is not cheapness but \emph{evidence before
commitment}.

The detection work needs no money at all: about a week on a departmental GPU
and 30\,GB of storage, against a dataset that is already public.

\subsection{If this build works}

There are two further steps, and each would be asked for separately, on the
evidence produced by the one before it. \textbf{We are not asking for either
now.}

\emph{Completing the aircraft} --- adding the video downlink and its ground
receivers, the RTK base station that makes positioning centimetre-accurate, the height
rangefinder and the coordination radio --- takes the build from
\rs{1{,}98{,}670} to \rs{2{,}45{,}993}. That is what turns a working aircraft
into a competition-legal one and makes the geolocation accuracy measurable
rather than merely designed for.

\emph{Building the other two} takes it to \rs{5{,}78{,}179} and is what buys
ten hectares in eight minutes. The replication is cheaper per aircraft than the
first, because the instruments, the charger, the safety-pilot transmitter and
the ground station are bought once and are already in this ask, and because
one RTK base station serves all three aircraft.

%-----------------------------------------------------------------------------
\section{Future work}

\textbf{Near term.} Measure what is currently calculated: detection accuracy,
motor thrust, setup time. Fly the full mission on one aircraft over its
three-and-a-half-hectare share, and settle whether the two-stage filter is
safe to rely on. Evaluate near-infrared sensing --- water
absorbs strongly and skin does not, so a person appears bright against dark
water. A camera without its infrared filter costs \rs{3{,}000}, and a public
dataset allows evaluation before purchase. Only the training data is
flood-specific; the same system suits earthquake, wildfire and avalanche
search.

\textbf{Then, on evidence.} Complete the aircraft, then build the second and
third and fly ten hectares in under eight minutes as a coordinated group. The
coordination, launch deconfliction and area-splitting software this needs is
already written and already tested in simulation with three aircraft; what it
has never had is three real ones.

\textbf{Longer term.} A capable aircraft deciding for itself within a small
power budget is a hard requirement in planetary aviation. \emph{Ingenuity} flew
72 times on Mars from a 1.8\,kg airframe in an atmosphere under one per cent of
Earth's density; \emph{Dragonfly} launches in 2028 for Titan. A command takes
minutes to arrive, so every decision is made on board. What transfers is what is
not flood-specific: onboard survey planning, perception run to an energy budget,
and safe behaviour without a link. Microcontroller-class processors, often
proposed for such missions, cannot resolve targets of this size, so a planetary
configuration needs a different operating point.

%-----------------------------------------------------------------------------
\section{Conclusion}

One aircraft, one operator, no network: people found in flooded ground and
reached with a relief kit, and coordinates a rescue coordinator can act on. The
design is finished. What is left is to build it and find out whether it works
--- and this is the cheapest honest way to find that out.

The department is never exposed to more than one step at a time. The first two
come to \rs{18{,}504} and settle whether these motors will lift this aircraft.
If they will not, we stop there and the department still has a thrust stand.
Nothing above \rs{26{,}000} is released until that measurement exists.

\textbf{And if the aircraft does what we expect it to, we would come back
twice}: once to complete it, and once to build the other two, so that ten
hectares can be searched in eight minutes with a reserve in hand --- which is
the one thing a single aircraft cannot do. Both requests would arrive with a
flown mission behind them rather than a calculation.

Floods come back every monsoon, and nothing of this kind exists at a size two
people can carry to the water's edge. We would like to begin, and we are asking
the department to release Phase 1.

%-----------------------------------------------------------------------------
% One DoSA approval letter per phase, generated from the same parts list as the
% schedule and the component tables above.
%============================================================================
% GENERATED FILE -- do not edit.
% Produced by tools/proposal/build_mvp_brief.py from
% hardware/bom/sourced_bom.py. Edit the BOM, not this file.
%============================================================================

\clearpage
\thispagestyle{empty}
\begingroup
\footnotesize
\setlength{\parskip}{0.26em}
% The brief sets loose table spacing for readability; a letter has to fit one
% page, so it is reset here rather than globally.
\setlength{\extrarowheight}{0pt}
\renewcommand{\arraystretch}{1.02}

% ---- letterhead, matching the brief's ---------------------------------------
\noindent
\raisebox{-0.5\height}{\includegraphics[height=11mm]{brainwave.pdf}}%
\hfill
\raisebox{-0.5\height}{\includegraphics[height=10mm]{thapar_logo.png}}

\vspace{2mm}\hrule\vspace{2.5mm}

\begin{center}
{\small\bfseries RescueSwarm --- An Autonomous Multi-UAV System for Post-Disaster Search, Localisation and Payload Delivery}
\end{center}

\vspace{1mm}\hrule\vspace{2.5mm}

\noindent\begin{minipage}[t]{0.62\textwidth}
To\\
The Dean, Student Affairs (DoSA)\\
\emph{Through:} President, Thapar Amateur Astronomers Society (TAAS)\\
Thapar Institute of Engineering \& Technology, Patiala
\end{minipage}\hfill
\begin{minipage}[t]{0.34\textwidth}
\raggedleft
Ref.: RS/PH-1\\[2.2mm]
Date: \rule{28mm}{0.4pt}
\end{minipage}

\vspace{1mm}
\noindent\textbf{Subject:} Fund Request under TAAS Society for the RescueSwarm
Project --- Phase 1 of 12, Establish thrust and mass measurement capability

\noindent Respected Ma'am,

We, a team of students of Thapar Institute of Engineering \& Technology, are
working under the guidance of Dr. Surbhi Sharma and Dr. Mamata Gulati on \textbf{RescueSwarm}, an
autonomous multi-UAV system for post-disaster search, survivor localisation and
payload delivery. The software and simulation framework is complete, and the
programme is now in hardware integration and testing, funded in 12 steps so
that each is approved only once the previous has produced a stated result.

\noindent\textbf{This request is Phase 1 of 12: Establish thrust and mass measurement capability.} This is the first step, and is gated only on this approval.
We request funding for the following components:

\noindent\makebox[\textwidth][c]{%
\scriptsize
\begin{tabular}{@{}>{\raggedright\arraybackslash}p{3.1cm}>{\raggedright\arraybackslash}p{6.9cm}rr@{}}
\toprule
\textbf{Component} & \textbf{Model} & \textbf{Qty} & \textbf{Cost (INR)} \\
\midrule
Fasteners & Hex socket-head assortment M3/M4/M5, SS304 & 1 & 1{,}140 \\
Load cell & REES52 20 kg load cell + HX711 breakout & 1 & 854 \\
Threadlocker & Loctite 243, 50 ml & 1 & 664 \\
Thrust-stand mast & SURYAKUSH metal 2.1 m tripod & 1 & 373 \\
Throttle source & Servo tester / ESC consistency tester, PPM & 1 & 111 \\
Load-cell amplifier & 7Semi HX711 24-bit ADC & 1 & 104 \\
\midrule
\multicolumn{3}{@{}r}{\textbf{Total}} & \textbf{3{,}246} \\
\bottomrule
\end{tabular}}

\noindent We kindly request approval and financial support of
\textbf{\rs{3{,}246}} under TAAS Society for the procurement of the components
above. This is the parts cost of this phase; the programme schedule additionally
carries tax where it is still owed, and a 15\,\% contingency.

\noindent\textbf{Team Members}\quad\emph{Mentors:} Dr. Surbhi Sharma and Dr. Mamata Gulati

\noindent\makebox[\textwidth][c]{%
\scriptsize
\begin{tabular}{@{}>{\raggedright\arraybackslash}p{3.0cm}>{\raggedright\arraybackslash}p{6.1cm}ll@{}}
\toprule
\textbf{Name} & \textbf{Department} & \textbf{Batch} & \textbf{Roll No.} \\
\midrule
Swastik Kumar & Electronics \& Communication Engineering & 2024--2028 & 1024060150 \\
Samuel Masih & Computer Science \& Engineering & 2024--2028 & 1024170044 \\
Prabhjot Singh & Robotics and Artificial Intelligence Engg. & 2024--2028 & 1024230021 \\
Manan Kapoor & Computer Science \& Engineering & 2024--2028 & 1024030467 \\
Manav Pansari & Electronics \& Communication Engineering & 2024--2028 & 1024060151 \\
Ravi Kant Raja & Robotics and Artificial Intelligence Engg. & 2024--2028 & 1024230032 \\
Neil Kishore & Mechanical Engineering & 2024--2028 & 1024080017 \\
\bottomrule
\end{tabular}}

\vspace{1.5mm}
\noindent\textbf{Approvals \& Signatures}

\vspace{8mm}
\noindent
\begin{tabular}{@{}p{5.3cm}p{5.3cm}p{5.4cm}@{}}
\hrulefill & \hrulefill & \hrulefill \\
\textbf{Dr. Surbhi Sharma} & \textbf{Dr. Mamata Gulati} & \textbf{President, TAAS} \\
{\footnotesize Mentor} & {\footnotesize Mentor} &
{\footnotesize Thapar Amateur Astronomers Society} \\
\end{tabular}

\vspace{5.5mm}
\noindent
\begin{tabular}{@{}p{8.2cm}p{8.2cm}@{}}
\hrulefill & \hrulefill \\
\textbf{Swastik Kumar} & \textbf{Date} \\
{\footnotesize Department of Electronics \& Communication Engineering, Roll No. 1024060150} & \\
{\footnotesize Submitted on behalf of the team} & \\
\end{tabular}

\vspace{2.5mm}
\noindent \fbox{\begin{minipage}{0.968\textwidth}
\vspace{0.8mm}
{\footnotesize \textbf{For office use --- Office of the Dean, Student Affairs}}

\vspace{6mm}
\begin{tabular}{@{}p{4.6cm}p{5.5cm}p{5.3cm}@{}}
\hrulefill & \hrulefill & \hrulefill \\
{\footnotesize Amount sanctioned (INR)} & {\footnotesize Signature} &
{\footnotesize Date} \\
\end{tabular}
\vspace{0.8mm}
\end{minipage}}
\endgroup

\clearpage
\thispagestyle{empty}
\begingroup
\footnotesize
\setlength{\parskip}{0.26em}
% The brief sets loose table spacing for readability; a letter has to fit one
% page, so it is reset here rather than globally.
\setlength{\extrarowheight}{0pt}
\renewcommand{\arraystretch}{1.02}

% ---- letterhead, matching the brief's ---------------------------------------
\noindent
\raisebox{-0.5\height}{\includegraphics[height=11mm]{brainwave.pdf}}%
\hfill
\raisebox{-0.5\height}{\includegraphics[height=10mm]{thapar_logo.png}}

\vspace{2mm}\hrule\vspace{2.5mm}

\begin{center}
{\small\bfseries RescueSwarm --- An Autonomous Multi-UAV System for Post-Disaster Search, Localisation and Payload Delivery}
\end{center}

\vspace{1mm}\hrule\vspace{2.5mm}

\noindent\begin{minipage}[t]{0.62\textwidth}
To\\
The Dean, Student Affairs (DoSA)\\
\emph{Through:} President, Thapar Amateur Astronomers Society (TAAS)\\
Thapar Institute of Engineering \& Technology, Patiala
\end{minipage}\hfill
\begin{minipage}[t]{0.34\textwidth}
\raggedleft
Ref.: RS/PH-2\\[2.2mm]
Date: \rule{28mm}{0.4pt}
\end{minipage}

\vspace{1mm}
\noindent\textbf{Subject:} Fund Request under TAAS Society for the RescueSwarm
Project --- Phase 2 of 12, Measure one motor, propeller and speed controller on the stand

\noindent Respected Ma'am,

We, a team of students of Thapar Institute of Engineering \& Technology, are
working under the guidance of Dr. Surbhi Sharma and Dr. Mamata Gulati on \textbf{RescueSwarm}, an
autonomous multi-UAV system for post-disaster search, survivor localisation and
payload delivery. The software and simulation framework is complete, and the
programme is now in hardware integration and testing, funded in 12 steps so
that each is approved only once the previous has produced a stated result.

\noindent\textbf{This request is Phase 2 of 12: Measure one motor, propeller and speed controller on the stand.} It falls due only once the previous step has produced its stated result: \emph{instruments commissioned}.
We request funding for the following components:

\noindent\makebox[\textwidth][c]{%
\scriptsize
\begin{tabular}{@{}>{\raggedright\arraybackslash}p{3.1cm}>{\raggedright\arraybackslash}p{6.9cm}rr@{}}
\toprule
\textbf{Component} & \textbf{Model} & \textbf{Qty} & \textbf{Cost (INR)} \\
\midrule
Speed controllers & Darkmatter VISHNU 50 MK-III, 4-in-1, 50 A & 1 & 5{,}189 \\
Motors & Tarot TL96020, 5008, 340 KV & 1 & 4{,}703 \\
Propellers & 1855 counter-rotating carbon fibre, CW+CCW pair & 1 & 1{,}641 \\
\midrule
\multicolumn{3}{@{}r}{\textbf{Total}} & \textbf{11{,}533} \\
\bottomrule
\end{tabular}}

\noindent We kindly request approval and financial support of
\textbf{\rs{11{,}533}} under TAAS Society for the procurement of the components
above. This is the parts cost of this phase; the programme schedule additionally
carries tax where it is still owed, and a 15\,\% contingency.

\noindent\textbf{Team Members}\quad\emph{Mentors:} Dr. Surbhi Sharma and Dr. Mamata Gulati

\noindent\makebox[\textwidth][c]{%
\scriptsize
\begin{tabular}{@{}>{\raggedright\arraybackslash}p{3.0cm}>{\raggedright\arraybackslash}p{6.1cm}ll@{}}
\toprule
\textbf{Name} & \textbf{Department} & \textbf{Batch} & \textbf{Roll No.} \\
\midrule
Swastik Kumar & Electronics \& Communication Engineering & 2024--2028 & 1024060150 \\
Samuel Masih & Computer Science \& Engineering & 2024--2028 & 1024170044 \\
Prabhjot Singh & Robotics and Artificial Intelligence Engg. & 2024--2028 & 1024230021 \\
Manan Kapoor & Computer Science \& Engineering & 2024--2028 & 1024030467 \\
Manav Pansari & Electronics \& Communication Engineering & 2024--2028 & 1024060151 \\
Ravi Kant Raja & Robotics and Artificial Intelligence Engg. & 2024--2028 & 1024230032 \\
Neil Kishore & Mechanical Engineering & 2024--2028 & 1024080017 \\
\bottomrule
\end{tabular}}

\vspace{1.5mm}
\noindent\textbf{Approvals \& Signatures}

\vspace{8mm}
\noindent
\begin{tabular}{@{}p{5.3cm}p{5.3cm}p{5.4cm}@{}}
\hrulefill & \hrulefill & \hrulefill \\
\textbf{Dr. Surbhi Sharma} & \textbf{Dr. Mamata Gulati} & \textbf{President, TAAS} \\
{\footnotesize Mentor} & {\footnotesize Mentor} &
{\footnotesize Thapar Amateur Astronomers Society} \\
\end{tabular}

\vspace{5.5mm}
\noindent
\begin{tabular}{@{}p{8.2cm}p{8.2cm}@{}}
\hrulefill & \hrulefill \\
\textbf{Swastik Kumar} & \textbf{Date} \\
{\footnotesize Department of Electronics \& Communication Engineering, Roll No. 1024060150} & \\
{\footnotesize Submitted on behalf of the team} & \\
\end{tabular}

\vspace{2.5mm}
\noindent \fbox{\begin{minipage}{0.968\textwidth}
\vspace{0.8mm}
{\footnotesize \textbf{For office use --- Office of the Dean, Student Affairs}}

\vspace{6mm}
\begin{tabular}{@{}p{4.6cm}p{5.5cm}p{5.3cm}@{}}
\hrulefill & \hrulefill & \hrulefill \\
{\footnotesize Amount sanctioned (INR)} & {\footnotesize Signature} &
{\footnotesize Date} \\
\end{tabular}
\vspace{0.8mm}
\end{minipage}}
\endgroup

\clearpage
\thispagestyle{empty}
\begingroup
\footnotesize
\setlength{\parskip}{0.26em}
% The brief sets loose table spacing for readability; a letter has to fit one
% page, so it is reset here rather than globally.
\setlength{\extrarowheight}{0pt}
\renewcommand{\arraystretch}{1.02}

% ---- letterhead, matching the brief's ---------------------------------------
\noindent
\raisebox{-0.5\height}{\includegraphics[height=11mm]{brainwave.pdf}}%
\hfill
\raisebox{-0.5\height}{\includegraphics[height=10mm]{thapar_logo.png}}

\vspace{2mm}\hrule\vspace{2.5mm}

\begin{center}
{\small\bfseries RescueSwarm --- An Autonomous Multi-UAV System for Post-Disaster Search, Localisation and Payload Delivery}
\end{center}

\vspace{1mm}\hrule\vspace{2.5mm}

\noindent\begin{minipage}[t]{0.62\textwidth}
To\\
The Dean, Student Affairs (DoSA)\\
\emph{Through:} President, Thapar Amateur Astronomers Society (TAAS)\\
Thapar Institute of Engineering \& Technology, Patiala
\end{minipage}\hfill
\begin{minipage}[t]{0.34\textwidth}
\raggedleft
Ref.: RS/PH-3\\[2.2mm]
Date: \rule{28mm}{0.4pt}
\end{minipage}

\vspace{1mm}
\noindent\textbf{Subject:} Fund Request under TAAS Society for the RescueSwarm
Project --- Phase 3 of 12, Aircraft -- autopilot and control link

\noindent Respected Ma'am,

We, a team of students of Thapar Institute of Engineering \& Technology, are
working under the guidance of Dr. Surbhi Sharma and Dr. Mamata Gulati on \textbf{RescueSwarm}, an
autonomous multi-UAV system for post-disaster search, survivor localisation and
payload delivery. The software and simulation framework is complete, and the
programme is now in hardware integration and testing, funded in 12 steps so
that each is approved only once the previous has produced a stated result.

\noindent\textbf{This request is Phase 3 of 12: Aircraft -- autopilot and control link.} It falls due only once the previous step has produced its stated result: \emph{thrust verified}.
We request funding for the following components:

\noindent\makebox[\textwidth][c]{%
\scriptsize
\begin{tabular}{@{}>{\raggedright\arraybackslash}p{3.1cm}>{\raggedright\arraybackslash}p{6.9cm}rr@{}}
\toprule
\textbf{Component} & \textbf{Model} & \textbf{Qty} & \textbf{Cost (INR)} \\
\midrule
Flight controller & Holybro Pixhawk 6C Mini & 1 & 22{,}049 \\
Autopilot log card & SanDisk 32 GB microSDHC, class 10 & 1 & 1{,}549 \\
FC vibration mount & Glass-fibre anti-vibration shock absorber & 1 & 161 \\
\midrule
\multicolumn{3}{@{}r}{\textbf{Total}} & \textbf{23{,}759} \\
\bottomrule
\end{tabular}}

\noindent We kindly request approval and financial support of
\textbf{\rs{23{,}759}} under TAAS Society for the procurement of the components
above. This is the parts cost of this phase; the programme schedule additionally
carries tax where it is still owed, and a 15\,\% contingency.

\noindent\textbf{Team Members}\quad\emph{Mentors:} Dr. Surbhi Sharma and Dr. Mamata Gulati

\noindent\makebox[\textwidth][c]{%
\scriptsize
\begin{tabular}{@{}>{\raggedright\arraybackslash}p{3.0cm}>{\raggedright\arraybackslash}p{6.1cm}ll@{}}
\toprule
\textbf{Name} & \textbf{Department} & \textbf{Batch} & \textbf{Roll No.} \\
\midrule
Swastik Kumar & Electronics \& Communication Engineering & 2024--2028 & 1024060150 \\
Samuel Masih & Computer Science \& Engineering & 2024--2028 & 1024170044 \\
Prabhjot Singh & Robotics and Artificial Intelligence Engg. & 2024--2028 & 1024230021 \\
Manan Kapoor & Computer Science \& Engineering & 2024--2028 & 1024030467 \\
Manav Pansari & Electronics \& Communication Engineering & 2024--2028 & 1024060151 \\
Ravi Kant Raja & Robotics and Artificial Intelligence Engg. & 2024--2028 & 1024230032 \\
Neil Kishore & Mechanical Engineering & 2024--2028 & 1024080017 \\
\bottomrule
\end{tabular}}

\vspace{1.5mm}
\noindent\textbf{Approvals \& Signatures}

\vspace{8mm}
\noindent
\begin{tabular}{@{}p{5.3cm}p{5.3cm}p{5.4cm}@{}}
\hrulefill & \hrulefill & \hrulefill \\
\textbf{Dr. Surbhi Sharma} & \textbf{Dr. Mamata Gulati} & \textbf{President, TAAS} \\
{\footnotesize Mentor} & {\footnotesize Mentor} &
{\footnotesize Thapar Amateur Astronomers Society} \\
\end{tabular}

\vspace{5.5mm}
\noindent
\begin{tabular}{@{}p{8.2cm}p{8.2cm}@{}}
\hrulefill & \hrulefill \\
\textbf{Swastik Kumar} & \textbf{Date} \\
{\footnotesize Department of Electronics \& Communication Engineering, Roll No. 1024060150} & \\
{\footnotesize Submitted on behalf of the team} & \\
\end{tabular}

\vspace{2.5mm}
\noindent \fbox{\begin{minipage}{0.968\textwidth}
\vspace{0.8mm}
{\footnotesize \textbf{For office use --- Office of the Dean, Student Affairs}}

\vspace{6mm}
\begin{tabular}{@{}p{4.6cm}p{5.5cm}p{5.3cm}@{}}
\hrulefill & \hrulefill & \hrulefill \\
{\footnotesize Amount sanctioned (INR)} & {\footnotesize Signature} &
{\footnotesize Date} \\
\end{tabular}
\vspace{0.8mm}
\end{minipage}}
\endgroup

\clearpage
\thispagestyle{empty}
\begingroup
\footnotesize
\setlength{\parskip}{0.26em}
% The brief sets loose table spacing for readability; a letter has to fit one
% page, so it is reset here rather than globally.
\setlength{\extrarowheight}{0pt}
\renewcommand{\arraystretch}{1.02}

% ---- letterhead, matching the brief's ---------------------------------------
\noindent
\raisebox{-0.5\height}{\includegraphics[height=11mm]{brainwave.pdf}}%
\hfill
\raisebox{-0.5\height}{\includegraphics[height=10mm]{thapar_logo.png}}

\vspace{2mm}\hrule\vspace{2.5mm}

\begin{center}
{\small\bfseries RescueSwarm --- An Autonomous Multi-UAV System for Post-Disaster Search, Localisation and Payload Delivery}
\end{center}

\vspace{1mm}\hrule\vspace{2.5mm}

\noindent\begin{minipage}[t]{0.62\textwidth}
To\\
The Dean, Student Affairs (DoSA)\\
\emph{Through:} President, Thapar Amateur Astronomers Society (TAAS)\\
Thapar Institute of Engineering \& Technology, Patiala
\end{minipage}\hfill
\begin{minipage}[t]{0.34\textwidth}
\raggedleft
Ref.: RS/PH-4\\[2.2mm]
Date: \rule{28mm}{0.4pt}
\end{minipage}

\vspace{1mm}
\noindent\textbf{Subject:} Fund Request under TAAS Society for the RescueSwarm
Project --- Phase 4 of 12, Aircraft -- onboard computer and AI accelerator

\noindent Respected Ma'am,

We, a team of students of Thapar Institute of Engineering \& Technology, are
working under the guidance of Dr. Surbhi Sharma and Dr. Mamata Gulati on \textbf{RescueSwarm}, an
autonomous multi-UAV system for post-disaster search, survivor localisation and
payload delivery. The software and simulation framework is complete, and the
programme is now in hardware integration and testing, funded in 12 steps so
that each is approved only once the previous has produced a stated result.

\noindent\textbf{This request is Phase 4 of 12: Aircraft -- onboard computer and AI accelerator.} It falls due only once the previous step has produced its stated result: \emph{autopilot bench-tested}.
We request funding for the following components:

\noindent\makebox[\textwidth][c]{%
\scriptsize
\begin{tabular}{@{}>{\raggedright\arraybackslash}p{3.1cm}>{\raggedright\arraybackslash}p{6.9cm}rr@{}}
\toprule
\textbf{Component} & \textbf{Model} & \textbf{Qty} & \textbf{Cost (INR)} \\
\midrule
Companion computer & Raspberry Pi 5, 8 GB & 1 & 19{,}999 \\
AI accelerator & Raspberry Pi AI HAT+, 26 TOPS & 1 & 11{,}749 \\
Storage & SanDisk High Endurance 128 GB microSD & 1 & 3{,}989 \\
Compute cooling & Official Raspberry Pi 5 active cooler & 1 & 519 \\
\midrule
\multicolumn{3}{@{}r}{\textbf{Total}} & \textbf{36{,}256} \\
\bottomrule
\end{tabular}}

\noindent We kindly request approval and financial support of
\textbf{\rs{36{,}256}} under TAAS Society for the procurement of the components
above. This is the parts cost of this phase; the programme schedule additionally
carries tax where it is still owed, and a 15\,\% contingency.

\noindent\textbf{Team Members}\quad\emph{Mentors:} Dr. Surbhi Sharma and Dr. Mamata Gulati

\noindent\makebox[\textwidth][c]{%
\scriptsize
\begin{tabular}{@{}>{\raggedright\arraybackslash}p{3.0cm}>{\raggedright\arraybackslash}p{6.1cm}ll@{}}
\toprule
\textbf{Name} & \textbf{Department} & \textbf{Batch} & \textbf{Roll No.} \\
\midrule
Swastik Kumar & Electronics \& Communication Engineering & 2024--2028 & 1024060150 \\
Samuel Masih & Computer Science \& Engineering & 2024--2028 & 1024170044 \\
Prabhjot Singh & Robotics and Artificial Intelligence Engg. & 2024--2028 & 1024230021 \\
Manan Kapoor & Computer Science \& Engineering & 2024--2028 & 1024030467 \\
Manav Pansari & Electronics \& Communication Engineering & 2024--2028 & 1024060151 \\
Ravi Kant Raja & Robotics and Artificial Intelligence Engg. & 2024--2028 & 1024230032 \\
Neil Kishore & Mechanical Engineering & 2024--2028 & 1024080017 \\
\bottomrule
\end{tabular}}

\vspace{1.5mm}
\noindent\textbf{Approvals \& Signatures}

\vspace{8mm}
\noindent
\begin{tabular}{@{}p{5.3cm}p{5.3cm}p{5.4cm}@{}}
\hrulefill & \hrulefill & \hrulefill \\
\textbf{Dr. Surbhi Sharma} & \textbf{Dr. Mamata Gulati} & \textbf{President, TAAS} \\
{\footnotesize Mentor} & {\footnotesize Mentor} &
{\footnotesize Thapar Amateur Astronomers Society} \\
\end{tabular}

\vspace{5.5mm}
\noindent
\begin{tabular}{@{}p{8.2cm}p{8.2cm}@{}}
\hrulefill & \hrulefill \\
\textbf{Swastik Kumar} & \textbf{Date} \\
{\footnotesize Department of Electronics \& Communication Engineering, Roll No. 1024060150} & \\
{\footnotesize Submitted on behalf of the team} & \\
\end{tabular}

\vspace{2.5mm}
\noindent \fbox{\begin{minipage}{0.968\textwidth}
\vspace{0.8mm}
{\footnotesize \textbf{For office use --- Office of the Dean, Student Affairs}}

\vspace{6mm}
\begin{tabular}{@{}p{4.6cm}p{5.5cm}p{5.3cm}@{}}
\hrulefill & \hrulefill & \hrulefill \\
{\footnotesize Amount sanctioned (INR)} & {\footnotesize Signature} &
{\footnotesize Date} \\
\end{tabular}
\vspace{0.8mm}
\end{minipage}}
\endgroup

\clearpage
\thispagestyle{empty}
\begingroup
\footnotesize
\setlength{\parskip}{0.26em}
% The brief sets loose table spacing for readability; a letter has to fit one
% page, so it is reset here rather than globally.
\setlength{\extrarowheight}{0pt}
\renewcommand{\arraystretch}{1.02}

% ---- letterhead, matching the brief's ---------------------------------------
\noindent
\raisebox{-0.5\height}{\includegraphics[height=11mm]{brainwave.pdf}}%
\hfill
\raisebox{-0.5\height}{\includegraphics[height=10mm]{thapar_logo.png}}

\vspace{2mm}\hrule\vspace{2.5mm}

\begin{center}
{\small\bfseries RescueSwarm --- An Autonomous Multi-UAV System for Post-Disaster Search, Localisation and Payload Delivery}
\end{center}

\vspace{1mm}\hrule\vspace{2.5mm}

\noindent\begin{minipage}[t]{0.62\textwidth}
To\\
The Dean, Student Affairs (DoSA)\\
\emph{Through:} President, Thapar Amateur Astronomers Society (TAAS)\\
Thapar Institute of Engineering \& Technology, Patiala
\end{minipage}\hfill
\begin{minipage}[t]{0.34\textwidth}
\raggedleft
Ref.: RS/PH-5\\[2.2mm]
Date: \rule{28mm}{0.4pt}
\end{minipage}

\vspace{1mm}
\noindent\textbf{Subject:} Fund Request under TAAS Society for the RescueSwarm
Project --- Phase 5 of 12, Aircraft -- satellite positioning and compass

\noindent Respected Ma'am,

We, a team of students of Thapar Institute of Engineering \& Technology, are
working under the guidance of Dr. Surbhi Sharma and Dr. Mamata Gulati on \textbf{RescueSwarm}, an
autonomous multi-UAV system for post-disaster search, survivor localisation and
payload delivery. The software and simulation framework is complete, and the
programme is now in hardware integration and testing, funded in 12 steps so
that each is approved only once the previous has produced a stated result.

\noindent\textbf{This request is Phase 5 of 12: Aircraft -- satellite positioning and compass.} It falls due only once the previous step has produced its stated result: \emph{compute stack operating}.
We request funding for the following components:

\noindent\makebox[\textwidth][c]{%
\scriptsize
\begin{tabular}{@{}>{\raggedright\arraybackslash}p{3.1cm}>{\raggedright\arraybackslash}p{6.9cm}rr@{}}
\toprule
\textbf{Component} & \textbf{Model} & \textbf{Qty} & \textbf{Cost (INR)} \\
\midrule
GNSS receiver & Holybro Micro M9N GPS with case, NEO-M9N and compass & 1 & 6{,}939 \\
\midrule
\multicolumn{3}{@{}r}{\textbf{Total}} & \textbf{6{,}939} \\
\bottomrule
\end{tabular}}

\noindent We kindly request approval and financial support of
\textbf{\rs{6{,}939}} under TAAS Society for the procurement of the components
above. This is the parts cost of this phase; the programme schedule additionally
carries tax where it is still owed, and a 15\,\% contingency.

\noindent\textbf{Team Members}\quad\emph{Mentors:} Dr. Surbhi Sharma and Dr. Mamata Gulati

\noindent\makebox[\textwidth][c]{%
\scriptsize
\begin{tabular}{@{}>{\raggedright\arraybackslash}p{3.0cm}>{\raggedright\arraybackslash}p{6.1cm}ll@{}}
\toprule
\textbf{Name} & \textbf{Department} & \textbf{Batch} & \textbf{Roll No.} \\
\midrule
Swastik Kumar & Electronics \& Communication Engineering & 2024--2028 & 1024060150 \\
Samuel Masih & Computer Science \& Engineering & 2024--2028 & 1024170044 \\
Prabhjot Singh & Robotics and Artificial Intelligence Engg. & 2024--2028 & 1024230021 \\
Manan Kapoor & Computer Science \& Engineering & 2024--2028 & 1024030467 \\
Manav Pansari & Electronics \& Communication Engineering & 2024--2028 & 1024060151 \\
Ravi Kant Raja & Robotics and Artificial Intelligence Engg. & 2024--2028 & 1024230032 \\
Neil Kishore & Mechanical Engineering & 2024--2028 & 1024080017 \\
\bottomrule
\end{tabular}}

\vspace{1.5mm}
\noindent\textbf{Approvals \& Signatures}

\vspace{8mm}
\noindent
\begin{tabular}{@{}p{5.3cm}p{5.3cm}p{5.4cm}@{}}
\hrulefill & \hrulefill & \hrulefill \\
\textbf{Dr. Surbhi Sharma} & \textbf{Dr. Mamata Gulati} & \textbf{President, TAAS} \\
{\footnotesize Mentor} & {\footnotesize Mentor} &
{\footnotesize Thapar Amateur Astronomers Society} \\
\end{tabular}

\vspace{5.5mm}
\noindent
\begin{tabular}{@{}p{8.2cm}p{8.2cm}@{}}
\hrulefill & \hrulefill \\
\textbf{Swastik Kumar} & \textbf{Date} \\
{\footnotesize Department of Electronics \& Communication Engineering, Roll No. 1024060150} & \\
{\footnotesize Submitted on behalf of the team} & \\
\end{tabular}

\vspace{2.5mm}
\noindent \fbox{\begin{minipage}{0.968\textwidth}
\vspace{0.8mm}
{\footnotesize \textbf{For office use --- Office of the Dean, Student Affairs}}

\vspace{6mm}
\begin{tabular}{@{}p{4.6cm}p{5.5cm}p{5.3cm}@{}}
\hrulefill & \hrulefill & \hrulefill \\
{\footnotesize Amount sanctioned (INR)} & {\footnotesize Signature} &
{\footnotesize Date} \\
\end{tabular}
\vspace{0.8mm}
\end{minipage}}
\endgroup

\clearpage
\thispagestyle{empty}
\begingroup
\footnotesize
\setlength{\parskip}{0.26em}
% The brief sets loose table spacing for readability; a letter has to fit one
% page, so it is reset here rather than globally.
\setlength{\extrarowheight}{0pt}
\renewcommand{\arraystretch}{1.02}

% ---- letterhead, matching the brief's ---------------------------------------
\noindent
\raisebox{-0.5\height}{\includegraphics[height=11mm]{brainwave.pdf}}%
\hfill
\raisebox{-0.5\height}{\includegraphics[height=10mm]{thapar_logo.png}}

\vspace{2mm}\hrule\vspace{2.5mm}

\begin{center}
{\small\bfseries RescueSwarm --- An Autonomous Multi-UAV System for Post-Disaster Search, Localisation and Payload Delivery}
\end{center}

\vspace{1mm}\hrule\vspace{2.5mm}

\noindent\begin{minipage}[t]{0.62\textwidth}
To\\
The Dean, Student Affairs (DoSA)\\
\emph{Through:} President, Thapar Amateur Astronomers Society (TAAS)\\
Thapar Institute of Engineering \& Technology, Patiala
\end{minipage}\hfill
\begin{minipage}[t]{0.34\textwidth}
\raggedleft
Ref.: RS/PH-6\\[2.2mm]
Date: \rule{28mm}{0.4pt}
\end{minipage}

\vspace{1mm}
\noindent\textbf{Subject:} Fund Request under TAAS Society for the RescueSwarm
Project --- Phase 6 of 12, Aircraft -- camera and command link

\noindent Respected Ma'am,

We, a team of students of Thapar Institute of Engineering \& Technology, are
working under the guidance of Dr. Surbhi Sharma and Dr. Mamata Gulati on \textbf{RescueSwarm}, an
autonomous multi-UAV system for post-disaster search, survivor localisation and
payload delivery. The software and simulation framework is complete, and the
programme is now in hardware integration and testing, funded in 12 steps so
that each is approved only once the previous has produced a stated result.

\noindent\textbf{This request is Phase 6 of 12: Aircraft -- camera and command link.} It falls due only once the previous step has produced its stated result: \emph{position fix acquired}.
We request funding for the following components:

\noindent\makebox[\textwidth][c]{%
\scriptsize
\begin{tabular}{@{}>{\raggedright\arraybackslash}p{3.1cm}>{\raggedright\arraybackslash}p{6.9cm}rr@{}}
\toprule
\textbf{Component} & \textbf{Model} & \textbf{Qty} & \textbf{Cost (INR)} \\
\midrule
Camera + lens & Arducam 12.3 MP 1/2.3 in HQ module, 6 mm CS lens & 1 & 6{,}799 \\
Command receiver & ExpressLRS RX24T, 2.4 GHz & 1 & 1{,}439 \\
\midrule
\multicolumn{3}{@{}r}{\textbf{Total}} & \textbf{8{,}238} \\
\bottomrule
\end{tabular}}

\noindent We kindly request approval and financial support of
\textbf{\rs{8{,}238}} under TAAS Society for the procurement of the components
above. This is the parts cost of this phase; the programme schedule additionally
carries tax where it is still owed, and a 15\,\% contingency.

\noindent\textbf{Team Members}\quad\emph{Mentors:} Dr. Surbhi Sharma and Dr. Mamata Gulati

\noindent\makebox[\textwidth][c]{%
\scriptsize
\begin{tabular}{@{}>{\raggedright\arraybackslash}p{3.0cm}>{\raggedright\arraybackslash}p{6.1cm}ll@{}}
\toprule
\textbf{Name} & \textbf{Department} & \textbf{Batch} & \textbf{Roll No.} \\
\midrule
Swastik Kumar & Electronics \& Communication Engineering & 2024--2028 & 1024060150 \\
Samuel Masih & Computer Science \& Engineering & 2024--2028 & 1024170044 \\
Prabhjot Singh & Robotics and Artificial Intelligence Engg. & 2024--2028 & 1024230021 \\
Manan Kapoor & Computer Science \& Engineering & 2024--2028 & 1024030467 \\
Manav Pansari & Electronics \& Communication Engineering & 2024--2028 & 1024060151 \\
Ravi Kant Raja & Robotics and Artificial Intelligence Engg. & 2024--2028 & 1024230032 \\
Neil Kishore & Mechanical Engineering & 2024--2028 & 1024080017 \\
\bottomrule
\end{tabular}}

\vspace{1.5mm}
\noindent\textbf{Approvals \& Signatures}

\vspace{8mm}
\noindent
\begin{tabular}{@{}p{5.3cm}p{5.3cm}p{5.4cm}@{}}
\hrulefill & \hrulefill & \hrulefill \\
\textbf{Dr. Surbhi Sharma} & \textbf{Dr. Mamata Gulati} & \textbf{President, TAAS} \\
{\footnotesize Mentor} & {\footnotesize Mentor} &
{\footnotesize Thapar Amateur Astronomers Society} \\
\end{tabular}

\vspace{5.5mm}
\noindent
\begin{tabular}{@{}p{8.2cm}p{8.2cm}@{}}
\hrulefill & \hrulefill \\
\textbf{Swastik Kumar} & \textbf{Date} \\
{\footnotesize Department of Electronics \& Communication Engineering, Roll No. 1024060150} & \\
{\footnotesize Submitted on behalf of the team} & \\
\end{tabular}

\vspace{2.5mm}
\noindent \fbox{\begin{minipage}{0.968\textwidth}
\vspace{0.8mm}
{\footnotesize \textbf{For office use --- Office of the Dean, Student Affairs}}

\vspace{6mm}
\begin{tabular}{@{}p{4.6cm}p{5.5cm}p{5.3cm}@{}}
\hrulefill & \hrulefill & \hrulefill \\
{\footnotesize Amount sanctioned (INR)} & {\footnotesize Signature} &
{\footnotesize Date} \\
\end{tabular}
\vspace{0.8mm}
\end{minipage}}
\endgroup

\clearpage
\thispagestyle{empty}
\begingroup
\footnotesize
\setlength{\parskip}{0.26em}
% The brief sets loose table spacing for readability; a letter has to fit one
% page, so it is reset here rather than globally.
\setlength{\extrarowheight}{0pt}
\renewcommand{\arraystretch}{1.02}

% ---- letterhead, matching the brief's ---------------------------------------
\noindent
\raisebox{-0.5\height}{\includegraphics[height=11mm]{brainwave.pdf}}%
\hfill
\raisebox{-0.5\height}{\includegraphics[height=10mm]{thapar_logo.png}}

\vspace{2mm}\hrule\vspace{2.5mm}

\begin{center}
{\small\bfseries RescueSwarm --- An Autonomous Multi-UAV System for Post-Disaster Search, Localisation and Payload Delivery}
\end{center}

\vspace{1mm}\hrule\vspace{2.5mm}

\noindent\begin{minipage}[t]{0.62\textwidth}
To\\
The Dean, Student Affairs (DoSA)\\
\emph{Through:} President, Thapar Amateur Astronomers Society (TAAS)\\
Thapar Institute of Engineering \& Technology, Patiala
\end{minipage}\hfill
\begin{minipage}[t]{0.34\textwidth}
\raggedleft
Ref.: RS/PH-7\\[2.2mm]
Date: \rule{28mm}{0.4pt}
\end{minipage}

\vspace{1mm}
\noindent\textbf{Subject:} Fund Request under TAAS Society for the RescueSwarm
Project --- Phase 7 of 12, Aircraft -- airframe fabrication

\noindent Respected Ma'am,

We, a team of students of Thapar Institute of Engineering \& Technology, are
working under the guidance of Dr. Surbhi Sharma and Dr. Mamata Gulati on \textbf{RescueSwarm}, an
autonomous multi-UAV system for post-disaster search, survivor localisation and
payload delivery. The software and simulation framework is complete, and the
programme is now in hardware integration and testing, funded in 12 steps so
that each is approved only once the previous has produced a stated result.

\noindent\textbf{This request is Phase 7 of 12: Aircraft -- airframe fabrication.} It falls due only once the previous step has produced its stated result: \emph{avionics integrated}.
We request funding for the following components:

\noindent\makebox[\textwidth][c]{%
\scriptsize
\begin{tabular}{@{}>{\raggedright\arraybackslash}p{3.1cm}>{\raggedright\arraybackslash}p{6.9cm}rr@{}}
\toprule
\textbf{Component} & \textbf{Model} & \textbf{Qty} & \textbf{Cost (INR)} \\
\midrule
Arm tube & 3K carbon fibre, OD25 x ID23 x 1000 mm & 2 & 4{,}518 \\
Motor mounts & Tarot 25 mm motor mount, TL9602 & 4 & 4{,}476 \\
Arm clamps & T-bolt hose clamps 23-25 mm, SS304, 4 pcs & 1 & 3{,}012 \\
Frame plate stock & 3K carbon fibre sheet, 300 x 200 x 2 mm & 1 & 2{,}895 \\
Printed parts & Pro-Range PETG-CF filament, 1.75 mm, 1 kg & 1 & 949 \\
Landing gear & F450/F550 landing skid & 1 & 549 \\
Suspension springs & 5 mm stainless double torsion spring & 10 & 20 \\
\midrule
\multicolumn{3}{@{}r}{\textbf{Total}} & \textbf{16{,}419} \\
\bottomrule
\end{tabular}}

\noindent We kindly request approval and financial support of
\textbf{\rs{16{,}419}} under TAAS Society for the procurement of the components
above. This is the parts cost of this phase; the programme schedule additionally
carries tax where it is still owed, and a 15\,\% contingency.

\noindent\textbf{Team Members}\quad\emph{Mentors:} Dr. Surbhi Sharma and Dr. Mamata Gulati

\noindent\makebox[\textwidth][c]{%
\scriptsize
\begin{tabular}{@{}>{\raggedright\arraybackslash}p{3.0cm}>{\raggedright\arraybackslash}p{6.1cm}ll@{}}
\toprule
\textbf{Name} & \textbf{Department} & \textbf{Batch} & \textbf{Roll No.} \\
\midrule
Swastik Kumar & Electronics \& Communication Engineering & 2024--2028 & 1024060150 \\
Samuel Masih & Computer Science \& Engineering & 2024--2028 & 1024170044 \\
Prabhjot Singh & Robotics and Artificial Intelligence Engg. & 2024--2028 & 1024230021 \\
Manan Kapoor & Computer Science \& Engineering & 2024--2028 & 1024030467 \\
Manav Pansari & Electronics \& Communication Engineering & 2024--2028 & 1024060151 \\
Ravi Kant Raja & Robotics and Artificial Intelligence Engg. & 2024--2028 & 1024230032 \\
Neil Kishore & Mechanical Engineering & 2024--2028 & 1024080017 \\
\bottomrule
\end{tabular}}

\vspace{1.5mm}
\noindent\textbf{Approvals \& Signatures}

\vspace{8mm}
\noindent
\begin{tabular}{@{}p{5.3cm}p{5.3cm}p{5.4cm}@{}}
\hrulefill & \hrulefill & \hrulefill \\
\textbf{Dr. Surbhi Sharma} & \textbf{Dr. Mamata Gulati} & \textbf{President, TAAS} \\
{\footnotesize Mentor} & {\footnotesize Mentor} &
{\footnotesize Thapar Amateur Astronomers Society} \\
\end{tabular}

\vspace{5.5mm}
\noindent
\begin{tabular}{@{}p{8.2cm}p{8.2cm}@{}}
\hrulefill & \hrulefill \\
\textbf{Swastik Kumar} & \textbf{Date} \\
{\footnotesize Department of Electronics \& Communication Engineering, Roll No. 1024060150} & \\
{\footnotesize Submitted on behalf of the team} & \\
\end{tabular}

\vspace{2.5mm}
\noindent \fbox{\begin{minipage}{0.968\textwidth}
\vspace{0.8mm}
{\footnotesize \textbf{For office use --- Office of the Dean, Student Affairs}}

\vspace{6mm}
\begin{tabular}{@{}p{4.6cm}p{5.5cm}p{5.3cm}@{}}
\hrulefill & \hrulefill & \hrulefill \\
{\footnotesize Amount sanctioned (INR)} & {\footnotesize Signature} &
{\footnotesize Date} \\
\end{tabular}
\vspace{0.8mm}
\end{minipage}}
\endgroup

\clearpage
\thispagestyle{empty}
\begingroup
\footnotesize
\setlength{\parskip}{0.26em}
% The brief sets loose table spacing for readability; a letter has to fit one
% page, so it is reset here rather than globally.
\setlength{\extrarowheight}{0pt}
\renewcommand{\arraystretch}{1.02}

% ---- letterhead, matching the brief's ---------------------------------------
\noindent
\raisebox{-0.5\height}{\includegraphics[height=11mm]{brainwave.pdf}}%
\hfill
\raisebox{-0.5\height}{\includegraphics[height=10mm]{thapar_logo.png}}

\vspace{2mm}\hrule\vspace{2.5mm}

\begin{center}
{\small\bfseries RescueSwarm --- An Autonomous Multi-UAV System for Post-Disaster Search, Localisation and Payload Delivery}
\end{center}

\vspace{1mm}\hrule\vspace{2.5mm}

\noindent\begin{minipage}[t]{0.62\textwidth}
To\\
The Dean, Student Affairs (DoSA)\\
\emph{Through:} President, Thapar Amateur Astronomers Society (TAAS)\\
Thapar Institute of Engineering \& Technology, Patiala
\end{minipage}\hfill
\begin{minipage}[t]{0.34\textwidth}
\raggedleft
Ref.: RS/PH-8\\[2.2mm]
Date: \rule{28mm}{0.4pt}
\end{minipage}

\vspace{1mm}
\noindent\textbf{Subject:} Fund Request under TAAS Society for the RescueSwarm
Project --- Phase 8 of 12, Aircraft -- battery pack and power distribution

\noindent Respected Ma'am,

We, a team of students of Thapar Institute of Engineering \& Technology, are
working under the guidance of Dr. Surbhi Sharma and Dr. Mamata Gulati on \textbf{RescueSwarm}, an
autonomous multi-UAV system for post-disaster search, survivor localisation and
payload delivery. The software and simulation framework is complete, and the
programme is now in hardware integration and testing, funded in 12 steps so
that each is approved only once the previous has produced a stated result.

\noindent\textbf{This request is Phase 8 of 12: Aircraft -- battery pack and power distribution.} It falls due only once the previous step has produced its stated result: \emph{airframe fabricated}.
We request funding for the following components:

\noindent\makebox[\textwidth][c]{%
\scriptsize
\begin{tabular}{@{}>{\raggedright\arraybackslash}p{3.1cm}>{\raggedright\arraybackslash}p{6.9cm}rr@{}}
\toprule
\textbf{Component} & \textbf{Model} & \textbf{Qty} & \textbf{Cost (INR)} \\
\midrule
Battery pack & GenX 22.2V 6S 12000mAh 25C/50C LiPo, XT90 & 1 & 9{,}088 \\
Power module & Holybro PM07 & 1 & 5{,}249 \\
BEC, primary & 8 A UBEC, 5/6 V out, 2-6S in & 1 & 3{,}939 \\
Pack fusing & ANL fuse holder + 150 A ANL fuses & 1 & 3{,}000 \\
BEC, secondary & ReadytoSky 2-6S 5V/3A + 12V/3A switchable UBEC & 1 & 279 \\
Pack retention & 300 mm LiPo strap, reusable & 2 & 124 \\
\midrule
\multicolumn{3}{@{}r}{\textbf{Total}} & \textbf{21{,}679} \\
\bottomrule
\end{tabular}}

\noindent We kindly request approval and financial support of
\textbf{\rs{21{,}679}} under TAAS Society for the procurement of the components
above. This is the parts cost of this phase; the programme schedule additionally
carries tax where it is still owed, and a 15\,\% contingency.

\noindent\textbf{Team Members}\quad\emph{Mentors:} Dr. Surbhi Sharma and Dr. Mamata Gulati

\noindent\makebox[\textwidth][c]{%
\scriptsize
\begin{tabular}{@{}>{\raggedright\arraybackslash}p{3.0cm}>{\raggedright\arraybackslash}p{6.1cm}ll@{}}
\toprule
\textbf{Name} & \textbf{Department} & \textbf{Batch} & \textbf{Roll No.} \\
\midrule
Swastik Kumar & Electronics \& Communication Engineering & 2024--2028 & 1024060150 \\
Samuel Masih & Computer Science \& Engineering & 2024--2028 & 1024170044 \\
Prabhjot Singh & Robotics and Artificial Intelligence Engg. & 2024--2028 & 1024230021 \\
Manan Kapoor & Computer Science \& Engineering & 2024--2028 & 1024030467 \\
Manav Pansari & Electronics \& Communication Engineering & 2024--2028 & 1024060151 \\
Ravi Kant Raja & Robotics and Artificial Intelligence Engg. & 2024--2028 & 1024230032 \\
Neil Kishore & Mechanical Engineering & 2024--2028 & 1024080017 \\
\bottomrule
\end{tabular}}

\vspace{1.5mm}
\noindent\textbf{Approvals \& Signatures}

\vspace{8mm}
\noindent
\begin{tabular}{@{}p{5.3cm}p{5.3cm}p{5.4cm}@{}}
\hrulefill & \hrulefill & \hrulefill \\
\textbf{Dr. Surbhi Sharma} & \textbf{Dr. Mamata Gulati} & \textbf{President, TAAS} \\
{\footnotesize Mentor} & {\footnotesize Mentor} &
{\footnotesize Thapar Amateur Astronomers Society} \\
\end{tabular}

\vspace{5.5mm}
\noindent
\begin{tabular}{@{}p{8.2cm}p{8.2cm}@{}}
\hrulefill & \hrulefill \\
\textbf{Swastik Kumar} & \textbf{Date} \\
{\footnotesize Department of Electronics \& Communication Engineering, Roll No. 1024060150} & \\
{\footnotesize Submitted on behalf of the team} & \\
\end{tabular}

\vspace{2.5mm}
\noindent \fbox{\begin{minipage}{0.968\textwidth}
\vspace{0.8mm}
{\footnotesize \textbf{For office use --- Office of the Dean, Student Affairs}}

\vspace{6mm}
\begin{tabular}{@{}p{4.6cm}p{5.5cm}p{5.3cm}@{}}
\hrulefill & \hrulefill & \hrulefill \\
{\footnotesize Amount sanctioned (INR)} & {\footnotesize Signature} &
{\footnotesize Date} \\
\end{tabular}
\vspace{0.8mm}
\end{minipage}}
\endgroup

\clearpage
\thispagestyle{empty}
\begingroup
\footnotesize
\setlength{\parskip}{0.26em}
% The brief sets loose table spacing for readability; a letter has to fit one
% page, so it is reset here rather than globally.
\setlength{\extrarowheight}{0pt}
\renewcommand{\arraystretch}{1.02}

% ---- letterhead, matching the brief's ---------------------------------------
\noindent
\raisebox{-0.5\height}{\includegraphics[height=11mm]{brainwave.pdf}}%
\hfill
\raisebox{-0.5\height}{\includegraphics[height=10mm]{thapar_logo.png}}

\vspace{2mm}\hrule\vspace{2.5mm}

\begin{center}
{\small\bfseries RescueSwarm --- An Autonomous Multi-UAV System for Post-Disaster Search, Localisation and Payload Delivery}
\end{center}

\vspace{1mm}\hrule\vspace{2.5mm}

\noindent\begin{minipage}[t]{0.62\textwidth}
To\\
The Dean, Student Affairs (DoSA)\\
\emph{Through:} President, Thapar Amateur Astronomers Society (TAAS)\\
Thapar Institute of Engineering \& Technology, Patiala
\end{minipage}\hfill
\begin{minipage}[t]{0.34\textwidth}
\raggedleft
Ref.: RS/PH-9\\[2.2mm]
Date: \rule{28mm}{0.4pt}
\end{minipage}

\vspace{1mm}
\noindent\textbf{Subject:} Fund Request under TAAS Society for the RescueSwarm
Project --- Phase 9 of 12, Aircraft -- release servos and wiring

\noindent Respected Ma'am,

We, a team of students of Thapar Institute of Engineering \& Technology, are
working under the guidance of Dr. Surbhi Sharma and Dr. Mamata Gulati on \textbf{RescueSwarm}, an
autonomous multi-UAV system for post-disaster search, survivor localisation and
payload delivery. The software and simulation framework is complete, and the
programme is now in hardware integration and testing, funded in 12 steps so
that each is approved only once the previous has produced a stated result.

\noindent\textbf{This request is Phase 9 of 12: Aircraft -- release servos and wiring.} It falls due only once the previous step has produced its stated result: \emph{pack bench-discharged}.
We request funding for the following components:

\noindent\makebox[\textwidth][c]{%
\scriptsize
\begin{tabular}{@{}>{\raggedright\arraybackslash}p{3.1cm}>{\raggedright\arraybackslash}p{6.9cm}rr@{}}
\toprule
\textbf{Component} & \textbf{Model} & \textbf{Qty} & \textbf{Cost (INR)} \\
\midrule
Antenna feeders & SMA M-F bulkhead RG316, 150 mm + adapters & 1 & 1{,}139 \\
Power connectors & Amass XT90S anti-spark, M/F pair & 4 & 720 \\
Release servos & MG90S mini servo, 180 deg & 4 & 596 \\
Main leads & 10 AWG ultra-flexible silicone wire & 1 & 149 \\
Insulation & PVC heat-shrink sleeve, 150 mm & 1 & 47 \\
Signal connectors & JST ZH 6-pin, 200 mm leads & 1 & 23 \\
\midrule
\multicolumn{3}{@{}r}{\textbf{Total}} & \textbf{2{,}674} \\
\bottomrule
\end{tabular}}

\noindent We kindly request approval and financial support of
\textbf{\rs{2{,}674}} under TAAS Society for the procurement of the components
above. This is the parts cost of this phase; the programme schedule additionally
carries tax where it is still owed, and a 15\,\% contingency.

\noindent\textbf{Team Members}\quad\emph{Mentors:} Dr. Surbhi Sharma and Dr. Mamata Gulati

\noindent\makebox[\textwidth][c]{%
\scriptsize
\begin{tabular}{@{}>{\raggedright\arraybackslash}p{3.0cm}>{\raggedright\arraybackslash}p{6.1cm}ll@{}}
\toprule
\textbf{Name} & \textbf{Department} & \textbf{Batch} & \textbf{Roll No.} \\
\midrule
Swastik Kumar & Electronics \& Communication Engineering & 2024--2028 & 1024060150 \\
Samuel Masih & Computer Science \& Engineering & 2024--2028 & 1024170044 \\
Prabhjot Singh & Robotics and Artificial Intelligence Engg. & 2024--2028 & 1024230021 \\
Manan Kapoor & Computer Science \& Engineering & 2024--2028 & 1024030467 \\
Manav Pansari & Electronics \& Communication Engineering & 2024--2028 & 1024060151 \\
Ravi Kant Raja & Robotics and Artificial Intelligence Engg. & 2024--2028 & 1024230032 \\
Neil Kishore & Mechanical Engineering & 2024--2028 & 1024080017 \\
\bottomrule
\end{tabular}}

\vspace{1.5mm}
\noindent\textbf{Approvals \& Signatures}

\vspace{8mm}
\noindent
\begin{tabular}{@{}p{5.3cm}p{5.3cm}p{5.4cm}@{}}
\hrulefill & \hrulefill & \hrulefill \\
\textbf{Dr. Surbhi Sharma} & \textbf{Dr. Mamata Gulati} & \textbf{President, TAAS} \\
{\footnotesize Mentor} & {\footnotesize Mentor} &
{\footnotesize Thapar Amateur Astronomers Society} \\
\end{tabular}

\vspace{5.5mm}
\noindent
\begin{tabular}{@{}p{8.2cm}p{8.2cm}@{}}
\hrulefill & \hrulefill \\
\textbf{Swastik Kumar} & \textbf{Date} \\
{\footnotesize Department of Electronics \& Communication Engineering, Roll No. 1024060150} & \\
{\footnotesize Submitted on behalf of the team} & \\
\end{tabular}

\vspace{2.5mm}
\noindent \fbox{\begin{minipage}{0.968\textwidth}
\vspace{0.8mm}
{\footnotesize \textbf{For office use --- Office of the Dean, Student Affairs}}

\vspace{6mm}
\begin{tabular}{@{}p{4.6cm}p{5.5cm}p{5.3cm}@{}}
\hrulefill & \hrulefill & \hrulefill \\
{\footnotesize Amount sanctioned (INR)} & {\footnotesize Signature} &
{\footnotesize Date} \\
\end{tabular}
\vspace{0.8mm}
\end{minipage}}
\endgroup

\clearpage
\thispagestyle{empty}
\begingroup
\footnotesize
\setlength{\parskip}{0.26em}
% The brief sets loose table spacing for readability; a letter has to fit one
% page, so it is reset here rather than globally.
\setlength{\extrarowheight}{0pt}
\renewcommand{\arraystretch}{1.02}

% ---- letterhead, matching the brief's ---------------------------------------
\noindent
\raisebox{-0.5\height}{\includegraphics[height=11mm]{brainwave.pdf}}%
\hfill
\raisebox{-0.5\height}{\includegraphics[height=10mm]{thapar_logo.png}}

\vspace{2mm}\hrule\vspace{2.5mm}

\begin{center}
{\small\bfseries RescueSwarm --- An Autonomous Multi-UAV System for Post-Disaster Search, Localisation and Payload Delivery}
\end{center}

\vspace{1mm}\hrule\vspace{2.5mm}

\noindent\begin{minipage}[t]{0.62\textwidth}
To\\
The Dean, Student Affairs (DoSA)\\
\emph{Through:} President, Thapar Amateur Astronomers Society (TAAS)\\
Thapar Institute of Engineering \& Technology, Patiala
\end{minipage}\hfill
\begin{minipage}[t]{0.34\textwidth}
\raggedleft
Ref.: RS/PH-10\\[2.2mm]
Date: \rule{28mm}{0.4pt}
\end{minipage}

\vspace{1mm}
\noindent\textbf{Subject:} Fund Request under TAAS Society for the RescueSwarm
Project --- Phase 10 of 12, Aircraft -- propulsion completed

\noindent Respected Ma'am,

We, a team of students of Thapar Institute of Engineering \& Technology, are
working under the guidance of Dr. Surbhi Sharma and Dr. Mamata Gulati on \textbf{RescueSwarm}, an
autonomous multi-UAV system for post-disaster search, survivor localisation and
payload delivery. The software and simulation framework is complete, and the
programme is now in hardware integration and testing, funded in 12 steps so
that each is approved only once the previous has produced a stated result.

\noindent\textbf{This request is Phase 10 of 12: Aircraft -- propulsion completed.} It falls due only once the previous step has produced its stated result: \emph{drive train installed}.
We request funding for the following components:

\noindent\makebox[\textwidth][c]{%
\scriptsize
\begin{tabular}{@{}>{\raggedright\arraybackslash}p{3.1cm}>{\raggedright\arraybackslash}p{6.9cm}rr@{}}
\toprule
\textbf{Component} & \textbf{Model} & \textbf{Qty} & \textbf{Cost (INR)} \\
\midrule
Motors & Tarot TL96020, 5008, 340 KV & 3 & 14{,}109 \\
Propellers & 1855 counter-rotating carbon fibre, CW+CCW pair & 5 & 8{,}205 \\
\midrule
\multicolumn{3}{@{}r}{\textbf{Total}} & \textbf{22{,}314} \\
\bottomrule
\end{tabular}}

\noindent We kindly request approval and financial support of
\textbf{\rs{22{,}314}} under TAAS Society for the procurement of the components
above. This is the parts cost of this phase; the programme schedule additionally
carries tax where it is still owed, and a 15\,\% contingency.

\noindent\textbf{Team Members}\quad\emph{Mentors:} Dr. Surbhi Sharma and Dr. Mamata Gulati

\noindent\makebox[\textwidth][c]{%
\scriptsize
\begin{tabular}{@{}>{\raggedright\arraybackslash}p{3.0cm}>{\raggedright\arraybackslash}p{6.1cm}ll@{}}
\toprule
\textbf{Name} & \textbf{Department} & \textbf{Batch} & \textbf{Roll No.} \\
\midrule
Swastik Kumar & Electronics \& Communication Engineering & 2024--2028 & 1024060150 \\
Samuel Masih & Computer Science \& Engineering & 2024--2028 & 1024170044 \\
Prabhjot Singh & Robotics and Artificial Intelligence Engg. & 2024--2028 & 1024230021 \\
Manan Kapoor & Computer Science \& Engineering & 2024--2028 & 1024030467 \\
Manav Pansari & Electronics \& Communication Engineering & 2024--2028 & 1024060151 \\
Ravi Kant Raja & Robotics and Artificial Intelligence Engg. & 2024--2028 & 1024230032 \\
Neil Kishore & Mechanical Engineering & 2024--2028 & 1024080017 \\
\bottomrule
\end{tabular}}

\vspace{1.5mm}
\noindent\textbf{Approvals \& Signatures}

\vspace{8mm}
\noindent
\begin{tabular}{@{}p{5.3cm}p{5.3cm}p{5.4cm}@{}}
\hrulefill & \hrulefill & \hrulefill \\
\textbf{Dr. Surbhi Sharma} & \textbf{Dr. Mamata Gulati} & \textbf{President, TAAS} \\
{\footnotesize Mentor} & {\footnotesize Mentor} &
{\footnotesize Thapar Amateur Astronomers Society} \\
\end{tabular}

\vspace{5.5mm}
\noindent
\begin{tabular}{@{}p{8.2cm}p{8.2cm}@{}}
\hrulefill & \hrulefill \\
\textbf{Swastik Kumar} & \textbf{Date} \\
{\footnotesize Department of Electronics \& Communication Engineering, Roll No. 1024060150} & \\
{\footnotesize Submitted on behalf of the team} & \\
\end{tabular}

\vspace{2.5mm}
\noindent \fbox{\begin{minipage}{0.968\textwidth}
\vspace{0.8mm}
{\footnotesize \textbf{For office use --- Office of the Dean, Student Affairs}}

\vspace{6mm}
\begin{tabular}{@{}p{4.6cm}p{5.5cm}p{5.3cm}@{}}
\hrulefill & \hrulefill & \hrulefill \\
{\footnotesize Amount sanctioned (INR)} & {\footnotesize Signature} &
{\footnotesize Date} \\
\end{tabular}
\vspace{0.8mm}
\end{minipage}}
\endgroup

\clearpage
\thispagestyle{empty}
\begingroup
\footnotesize
\setlength{\parskip}{0.26em}
% The brief sets loose table spacing for readability; a letter has to fit one
% page, so it is reset here rather than globally.
\setlength{\extrarowheight}{0pt}
\renewcommand{\arraystretch}{1.02}

% ---- letterhead, matching the brief's ---------------------------------------
\noindent
\raisebox{-0.5\height}{\includegraphics[height=11mm]{brainwave.pdf}}%
\hfill
\raisebox{-0.5\height}{\includegraphics[height=10mm]{thapar_logo.png}}

\vspace{2mm}\hrule\vspace{2.5mm}

\begin{center}
{\small\bfseries RescueSwarm --- An Autonomous Multi-UAV System for Post-Disaster Search, Localisation and Payload Delivery}
\end{center}

\vspace{1mm}\hrule\vspace{2.5mm}

\noindent\begin{minipage}[t]{0.62\textwidth}
To\\
The Dean, Student Affairs (DoSA)\\
\emph{Through:} President, Thapar Amateur Astronomers Society (TAAS)\\
Thapar Institute of Engineering \& Technology, Patiala
\end{minipage}\hfill
\begin{minipage}[t]{0.34\textwidth}
\raggedleft
Ref.: RS/PH-11\\[2.2mm]
Date: \rule{28mm}{0.4pt}
\end{minipage}

\vspace{1mm}
\noindent\textbf{Subject:} Fund Request under TAAS Society for the RescueSwarm
Project --- Phase 11 of 12, Charging and manual override

\noindent Respected Ma'am,

We, a team of students of Thapar Institute of Engineering \& Technology, are
working under the guidance of Dr. Surbhi Sharma and Dr. Mamata Gulati on \textbf{RescueSwarm}, an
autonomous multi-UAV system for post-disaster search, survivor localisation and
payload delivery. The software and simulation framework is complete, and the
programme is now in hardware integration and testing, funded in 12 steps so
that each is approved only once the previous has produced a stated result.

\noindent\textbf{This request is Phase 11 of 12: Charging and manual override.} It falls due only once the previous step has produced its stated result: \emph{aircraft assembled and weighed}.
We request funding for the following components:

\noindent\makebox[\textwidth][c]{%
\scriptsize
\begin{tabular}{@{}>{\raggedright\arraybackslash}p{3.1cm}>{\raggedright\arraybackslash}p{6.9cm}rr@{}}
\toprule
\textbf{Component} & \textbf{Model} & \textbf{Qty} & \textbf{Cost (INR)} \\
\midrule
Safety-pilot transmitter & RadioMaster TX12 MKII ELRS + RP1 receiver & 1 & 16{,}519 \\
Battery charger & ToolKitRC M6D, 500 W, 25 A, dual & 1 & 7{,}349 \\
\midrule
\multicolumn{3}{@{}r}{\textbf{Total}} & \textbf{23{,}868} \\
\bottomrule
\end{tabular}}

\noindent We kindly request approval and financial support of
\textbf{\rs{23{,}868}} under TAAS Society for the procurement of the components
above. This is the parts cost of this phase; the programme schedule additionally
carries tax where it is still owed, and a 15\,\% contingency.

\noindent\textbf{Team Members}\quad\emph{Mentors:} Dr. Surbhi Sharma and Dr. Mamata Gulati

\noindent\makebox[\textwidth][c]{%
\scriptsize
\begin{tabular}{@{}>{\raggedright\arraybackslash}p{3.0cm}>{\raggedright\arraybackslash}p{6.1cm}ll@{}}
\toprule
\textbf{Name} & \textbf{Department} & \textbf{Batch} & \textbf{Roll No.} \\
\midrule
Swastik Kumar & Electronics \& Communication Engineering & 2024--2028 & 1024060150 \\
Samuel Masih & Computer Science \& Engineering & 2024--2028 & 1024170044 \\
Prabhjot Singh & Robotics and Artificial Intelligence Engg. & 2024--2028 & 1024230021 \\
Manan Kapoor & Computer Science \& Engineering & 2024--2028 & 1024030467 \\
Manav Pansari & Electronics \& Communication Engineering & 2024--2028 & 1024060151 \\
Ravi Kant Raja & Robotics and Artificial Intelligence Engg. & 2024--2028 & 1024230032 \\
Neil Kishore & Mechanical Engineering & 2024--2028 & 1024080017 \\
\bottomrule
\end{tabular}}

\vspace{1.5mm}
\noindent\textbf{Approvals \& Signatures}

\vspace{8mm}
\noindent
\begin{tabular}{@{}p{5.3cm}p{5.3cm}p{5.4cm}@{}}
\hrulefill & \hrulefill & \hrulefill \\
\textbf{Dr. Surbhi Sharma} & \textbf{Dr. Mamata Gulati} & \textbf{President, TAAS} \\
{\footnotesize Mentor} & {\footnotesize Mentor} &
{\footnotesize Thapar Amateur Astronomers Society} \\
\end{tabular}

\vspace{5.5mm}
\noindent
\begin{tabular}{@{}p{8.2cm}p{8.2cm}@{}}
\hrulefill & \hrulefill \\
\textbf{Swastik Kumar} & \textbf{Date} \\
{\footnotesize Department of Electronics \& Communication Engineering, Roll No. 1024060150} & \\
{\footnotesize Submitted on behalf of the team} & \\
\end{tabular}

\vspace{2.5mm}
\noindent \fbox{\begin{minipage}{0.968\textwidth}
\vspace{0.8mm}
{\footnotesize \textbf{For office use --- Office of the Dean, Student Affairs}}

\vspace{6mm}
\begin{tabular}{@{}p{4.6cm}p{5.5cm}p{5.3cm}@{}}
\hrulefill & \hrulefill & \hrulefill \\
{\footnotesize Amount sanctioned (INR)} & {\footnotesize Signature} &
{\footnotesize Date} \\
\end{tabular}
\vspace{0.8mm}
\end{minipage}}
\endgroup

\clearpage
\thispagestyle{empty}
\begingroup
\footnotesize
\setlength{\parskip}{0.26em}
% The brief sets loose table spacing for readability; a letter has to fit one
% page, so it is reset here rather than globally.
\setlength{\extrarowheight}{0pt}
\renewcommand{\arraystretch}{1.02}

% ---- letterhead, matching the brief's ---------------------------------------
\noindent
\raisebox{-0.5\height}{\includegraphics[height=11mm]{brainwave.pdf}}%
\hfill
\raisebox{-0.5\height}{\includegraphics[height=10mm]{thapar_logo.png}}

\vspace{2mm}\hrule\vspace{2.5mm}

\begin{center}
{\small\bfseries RescueSwarm --- An Autonomous Multi-UAV System for Post-Disaster Search, Localisation and Payload Delivery}
\end{center}

\vspace{1mm}\hrule\vspace{2.5mm}

\noindent\begin{minipage}[t]{0.62\textwidth}
To\\
The Dean, Student Affairs (DoSA)\\
\emph{Through:} President, Thapar Amateur Astronomers Society (TAAS)\\
Thapar Institute of Engineering \& Technology, Patiala
\end{minipage}\hfill
\begin{minipage}[t]{0.34\textwidth}
\raggedleft
Ref.: RS/PH-12\\[2.2mm]
Date: \rule{28mm}{0.4pt}
\end{minipage}

\vspace{1mm}
\noindent\textbf{Subject:} Fund Request under TAAS Society for the RescueSwarm
Project --- Phase 12 of 12, Field consumables and assembly materials

\noindent Respected Ma'am,

We, a team of students of Thapar Institute of Engineering \& Technology, are
working under the guidance of Dr. Surbhi Sharma and Dr. Mamata Gulati on \textbf{RescueSwarm}, an
autonomous multi-UAV system for post-disaster search, survivor localisation and
payload delivery. The software and simulation framework is complete, and the
programme is now in hardware integration and testing, funded in 12 steps so
that each is approved only once the previous has produced a stated result.

\noindent\textbf{This request is Phase 12 of 12: Field consumables and assembly materials.} It falls due only once the previous step has produced its stated result: \emph{charging verified}.
We request funding for the following components:

\noindent\makebox[\textwidth][c]{%
\scriptsize
\begin{tabular}{@{}>{\raggedright\arraybackslash}p{3.1cm}>{\raggedright\arraybackslash}p{6.9cm}rr@{}}
\toprule
\textbf{Component} & \textbf{Model} & \textbf{Qty} & \textbf{Cost (INR)} \\
\midrule
Consumables & Fasteners, tape, connectors, filament & 1 & 5{,}000 \\
Mounting tape & 3M VHB 5952, 6 mm x 8.2 m & 1 & 560 \\
Heat-shrink kit & 328 pcs assorted 2:1 heat-shrink & 1 & 159 \\
Hook and loop & Self-adhesive hook and loop, 25 mm x 2 m & 1 & 145 \\
Cable management & Nylon zip ties, 350 mm, 100 pcs & 1 & 135 \\
\midrule
\multicolumn{3}{@{}r}{\textbf{Total}} & \textbf{5{,}999} \\
\bottomrule
\end{tabular}}

\noindent We kindly request approval and financial support of
\textbf{\rs{5{,}999}} under TAAS Society for the procurement of the components
above. This is the parts cost of this phase; the programme schedule additionally
carries tax where it is still owed, and a 15\,\% contingency.

\noindent\textbf{Team Members}\quad\emph{Mentors:} Dr. Surbhi Sharma and Dr. Mamata Gulati

\noindent\makebox[\textwidth][c]{%
\scriptsize
\begin{tabular}{@{}>{\raggedright\arraybackslash}p{3.0cm}>{\raggedright\arraybackslash}p{6.1cm}ll@{}}
\toprule
\textbf{Name} & \textbf{Department} & \textbf{Batch} & \textbf{Roll No.} \\
\midrule
Swastik Kumar & Electronics \& Communication Engineering & 2024--2028 & 1024060150 \\
Samuel Masih & Computer Science \& Engineering & 2024--2028 & 1024170044 \\
Prabhjot Singh & Robotics and Artificial Intelligence Engg. & 2024--2028 & 1024230021 \\
Manan Kapoor & Computer Science \& Engineering & 2024--2028 & 1024030467 \\
Manav Pansari & Electronics \& Communication Engineering & 2024--2028 & 1024060151 \\
Ravi Kant Raja & Robotics and Artificial Intelligence Engg. & 2024--2028 & 1024230032 \\
Neil Kishore & Mechanical Engineering & 2024--2028 & 1024080017 \\
\bottomrule
\end{tabular}}

\vspace{1.5mm}
\noindent\textbf{Approvals \& Signatures}

\vspace{8mm}
\noindent
\begin{tabular}{@{}p{5.3cm}p{5.3cm}p{5.4cm}@{}}
\hrulefill & \hrulefill & \hrulefill \\
\textbf{Dr. Surbhi Sharma} & \textbf{Dr. Mamata Gulati} & \textbf{President, TAAS} \\
{\footnotesize Mentor} & {\footnotesize Mentor} &
{\footnotesize Thapar Amateur Astronomers Society} \\
\end{tabular}

\vspace{5.5mm}
\noindent
\begin{tabular}{@{}p{8.2cm}p{8.2cm}@{}}
\hrulefill & \hrulefill \\
\textbf{Swastik Kumar} & \textbf{Date} \\
{\footnotesize Department of Electronics \& Communication Engineering, Roll No. 1024060150} & \\
{\footnotesize Submitted on behalf of the team} & \\
\end{tabular}

\vspace{2.5mm}
\noindent \fbox{\begin{minipage}{0.968\textwidth}
\vspace{0.8mm}
{\footnotesize \textbf{For office use --- Office of the Dean, Student Affairs}}

\vspace{6mm}
\begin{tabular}{@{}p{4.6cm}p{5.5cm}p{5.3cm}@{}}
\hrulefill & \hrulefill & \hrulefill \\
{\footnotesize Amount sanctioned (INR)} & {\footnotesize Signature} &
{\footnotesize Date} \\
\end{tabular}
\vspace{0.8mm}
\end{minipage}}
\endgroup


\end{document}
